% ================================================================
% MEANING GAMES: CHAPTERS 1--7 (EXPANDED)
%
% A Reply to Wittgenstein and the Game Theorists
%
% No preamble --- included by the main document.
% Assumes: \rs, \emerge (\kappa), \compose (\circ),
%          \void (\varepsilon), \IS (\mathcal I), \R, \N
%          Theorem environments: theorem, lemma, proposition,
%          corollary, definition, example, axiom_def, game,
%          remark, interpretation, aside
%          Lean listing environment via lstlisting
% ================================================================



%==========================================================================
\part{The Structure of Meaning Games}
%==========================================================================


% ================================================================
% CHAPTER 1: WITTGENSTEIN'S LANGUAGE GAMES
% ================================================================
\chapter{Wittgenstein's Language Games}
\label{ch:wittgenstein}

\epigraph{``The meaning of a word is its use in the language.''}
{---Ludwig Wittgenstein, \emph{Philosophical Investigations} \S43}


\section{The Philosophical Problem}

Wittgenstein opens the \emph{Philosophical Investigations} with a quotation
from Augustine's \emph{Confessions}:
\begin{quote}
``When they (my elders) named some object, and accordingly moved towards
something, I saw this and I grasped that the thing was called by the sound
they uttered when they meant to point it out.'' (\emph{PI} \S1)
\end{quote}
This picture of language---each word names an object, each sentence a
combination of names---is what Wittgenstein sets out to destroy. The
Augustinian picture assumes that \emph{meaning is reference}: a word means
the object it points to. Wittgenstein replaces this with: \emph{meaning is
use}.

But Wittgenstein never provides a mathematical framework for ``use.'' The
\emph{Investigations} are deliberately anti-systematic, proceeding by example
and counterexample, by thought experiment and therapeutic dissolution.
``Philosophy,'' Wittgenstein writes, ``is a battle against the bewitchment of
our intelligence by means of our language'' (\emph{PI} \S109).

We propose that \textbf{Idea Diffusion Theory provides the mathematics that
Wittgenstein's philosophy demands but refuses to supply}. The resonance
function $\rs$ is the formal correlate of ``use'': the meaning of an
utterance $a$ is not what $a$ \emph{refers to}, but how $a$ resonates---with
other utterances, with contexts, with observers. Two utterances have the same
meaning if and only if they have the same resonance profile:
\[
  a \sim b \iff \forall c \in \IS,\; \rs(a, c) = \rs(b, c)
  \quad\text{and}\quad \rs(c, a) = \rs(c, b).
\]
This is ``meaning as use,'' stated with the precision that Wittgenstein
himself would have both demanded and distrusted.


\subsection{The Augustinian Picture, Formalized}

The Augustinian picture of language corresponds to a degenerate ideatic space.

\begin{definition}[Augustinian Ideatic Space]
\label{def:augustinian}
An ideatic space $(\IS, \compose, \void, \rs)$ is \emph{Augustinian} if:
\begin{enumerate}
\item There exists a set $O$ of ``objects'' and a bijection $\nu : \IS \setminus \{\void\} \to O$ (naming),
\item $\rs(a, b) = \begin{cases} w(a) & \text{if } \nu(a) = \nu(b), \\ 0 & \text{otherwise}.\end{cases}$
\end{enumerate}
In an Augustinian space, resonance reduces to identity of reference:
two words resonate if and only if they name the same object.
\end{definition}

\begin{proposition}[Augustinian Spaces Have Zero Emergence]
\label{prop:augustinian-zero-emergence}
In any Augustinian ideatic space, $\emerge(a, b, c) = 0$ for all $a, b, c$.
\end{proposition}

\begin{proof}
In an Augustinian space, $\rs(x, c) = w(x) \cdot \mathbf{1}[\nu(x) = \nu(c)]$
for all $x, c$. The emergence is
\[
\emerge(a,b,c) = \rs(a \compose b, c) - \rs(a, c) - \rs(b, c).
\]
If $a \compose b$ does not name $\nu(c)$, then $\rs(a \compose b, c) = 0$. But
since $a$ and $b$ individually either do or do not name $\nu(c)$, and naming
is a function, we get $\rs(a,c) + \rs(b,c) = 0$ as well (at most one can name
$\nu(c)$, but in the Augustinian model composition produces a single name).
In all cases the emergence vanishes.
\end{proof}

\begin{interpretation}
The Augustinian picture of language is the $\emerge \equiv 0$ special case
of IDT. Wittgenstein's \emph{Investigations} can be read as a sustained
argument that $\emerge \not\equiv 0$---that the composition of words creates
meaning that cannot be decomposed into the meanings of the parts. The
emergence term is the mathematics of Wittgenstein's protest.
\end{interpretation}

\begin{interpretation}[Connection to Lewis's \emph{Convention}]
David Lewis's \emph{Convention} (1969) formalized language as a solution to
coordination games. In Lewis's framework, a convention is a regularity in
behavior that solves a recurrent coordination problem: everyone drives on
the right because everyone else does. Lewis's account is purely
\emph{exchange-based}: agents coordinate on signals whose value is fixed
by the coordination equilibrium. There is no emergence in Lewis---the
``meaning'' of driving on the right is exhausted by the fact that others
do so.

IDT reveals what Lewis's account misses: \emph{creative coordination}.
When a community converges on using ``cool'' to mean ``socially desirable''
(rather than ``low temperature''), the new meaning is not merely a
coordination equilibrium---it is an \emph{emergent} resonance that arises
from the interaction of the word with the community's evolving context.
The convention creates meaning ($\emerge > 0$) that no individual agent
intended. Lewis gives us coordination; IDT gives us coordination
\emph{plus creation}.
\end{interpretation}

\begin{remark}[Brandom's Normative Pragmatics]
Robert Brandom's \emph{Making It Explicit} (1994) argues that meaning
arises from the \emph{normative} practices of giving and asking for
reasons. An utterance's meaning is constituted by the inferential
commitments and entitlements it carries. In IDT terms, Brandom's
``score-keeping'' model corresponds to tracking resonance profiles:
the ``score'' of an utterance $a$ in a social context is precisely
its resonance profile $\rho_a$. The normative dimension---that certain
inferences are \emph{correct}---corresponds to the sign of emergence:
positive emergence marks warranted inference, negative emergence marks
fallacy. Where Brandom relies on normative vocabulary (``commitment,''
``entitlement''), IDT provides a quantitative substrate.
\end{remark}

\begin{lstlisting}[caption={Lean 4: Augustinian ideatic space has zero emergence}]
/-- An ideatic space where resonance reduces to name-matching
    has identically zero emergence. -/
theorem augustinian_zero_emergence
    {I : Type*} [S : IdeaticSpace I]
    (naming : I → ℕ)
    (h_rs : ∀ a b : I, rs a b = if naming a = naming b
      then rs a a else 0) :
    ∀ a b c : I, emergence a b c = 0 := by
  intro a b c
  unfold emergence
  rw [h_rs (compose a b) c, h_rs a c, h_rs b c]
  split_ifs <;> linarith
\end{lstlisting}

\begin{example}[The Augustinian Limit: Machine Translation]
Early machine translation systems (1950s--1970s) operated in an essentially
Augustinian framework: each word had a fixed ``meaning'' (its translation
in the target language), and translation was word-by-word substitution.
These systems failed spectacularly on idiomatic language, metaphor, and
context-dependent meaning---precisely because they assumed $\emerge = 0$.
Modern neural machine translation (Vaswani et~al.\ 2017) succeeds because
it implicitly models emergence: the ``attention mechanism'' captures how
words interact nonlinearly with their context to produce meaning. The
Transformer architecture is, in this sense, an empirical discovery of
emergence.
\end{example}

\begin{remark}[The Tractatus vs.\ The Investigations]
Wittgenstein's own philosophical development traces the arc from
Augustinian ($\emerge = 0$) to emergent ($\emerge \neq 0$) theories
of meaning. The \emph{Tractatus Logico-Philosophicus} (1921) posits a
``picture theory'' of meaning in which propositions picture facts
through a shared logical form---a theory with no room for emergence,
since the meaning of a proposition is determined entirely by its
logical structure and the objects it names.

The \emph{Philosophical Investigations} (1953) repudiates the
\emph{Tractatus} precisely on this point: meaning is not determined
by logical structure alone but by \emph{use in a form of life}. Use
involves context, history, social practice, and all the nonlinear
interactions that IDT captures through the emergence term. The
\emph{Tractatus} is the $\emerge = 0$ theory; the \emph{Investigations}
is the insistence that $\emerge \neq 0$. IDT provides the mathematics
for both, and shows that the former is a limiting case of the latter.
\end{remark}


\section{Language Games as Meaning Games}

Wittgenstein introduces the concept of \emph{Sprachspiel} (language game) in
\S7 of the \emph{Investigations}:
\begin{quote}
``I shall also call the whole, consisting of language and the actions into
which it is woven, the `language-game.'\,'' (\emph{PI} \S7)
\end{quote}
The crucial word is ``woven'': language is not applied \emph{to} an activity
from outside; it is woven \emph{into} the activity, inseparable from it. In
\S23 Wittgenstein lists examples of language games: giving orders, describing
objects, reporting events, speculating about events, forming and testing
hypotheses, presenting results of an experiment in tables and diagrams,
making up stories, singing, guessing riddles, making jokes, translating,
asking, thanking, cursing, greeting, and praying.

\begin{definition}[Language Game]
\label{def:language-game-formal}
A \emph{language game} in the sense of IDT is a tuple
$\Gamma = (\IS, \compose, \void, \rs, a, b)$ where:
\begin{itemize}
\item $(\IS, \compose, \void, \rs)$ is an ideatic space,
\item $a \in \IS$ is the \emph{intent} (what the sender aims to convey),
\item $b \in \IS$ is the \emph{form of life} (the receiver's context).
\end{itemize}
The sender's strategy set is $\IS$ itself: any utterance can be chosen as
signal. The payoff from signal $s$ is
\[
  u(s; a, b) := \rs(s \compose b, a).
\]
This measures how well the received idea $s \compose b$ resonates with the
sender's intent.
\end{definition}

\begin{theorem}[Universality of the Language Game]
\label{thm:universality}
Every act of communication in an ideatic space---sending orders, telling
stories, making jokes, praying---is an instance of a language game
$\Gamma = (\IS, \compose, \void, \rs, a, b)$ with appropriate choice of
$a$ and $b$. The mathematical structure is universal; only the parameters
differ.
\end{theorem}

\begin{proof}
This is definitional but warrants careful analysis. The sender has some
intent $a$; the receiver has some context $b$; the sender chooses a signal
$s$; the received idea is $s \compose b$. The resonance
$\rs(s \compose b, a)$ measures success. Every act of communication
instantiates this pattern.

What differs is the ideatic space $\IS$ (which ideas are available), the
composition $\compose$ (how ideas combine), the resonance function $\rs$
(what counts as success), and the specific parameters $a$ and $b$.

The claim of universality is strong: it asserts that there is \emph{no}
communicative act that falls outside this framework. Ordering a slab,
telling a joke, making a scientific argument, praying, and cursing all
have the same \emph{mathematical form}---they differ only in the choice
of ideatic space and parameters. This universality is Wittgenstein's
insight that ``the concept `game' is a concept with blurred edges''
(\emph{PI} \S71), combined with the mathematical precision that
Wittgenstein refused: the edges are ``blurred'' because the parameters
vary continuously, not because the structure is absent.
\end{proof}

\begin{example}[Ordering a Slab]
Wittgenstein's builder's game (\emph{PI} \S2): a builder $A$ calls out
``Slab!'' and helper $B$ brings a slab. Here $a = \text{``need-slab''}$,
$b = \text{``builder-context''}$, $s = \text{``Slab!''}$. The language game
succeeds when $\rs(\text{``Slab!''} \compose \text{``builder-context''},
\text{``need-slab''}) > 0$: the composition of the word with the shared
form of life resonates with the intent.
\end{example}

\begin{example}[Courtroom Argument]
A defense attorney has intent $a = \text{``reasonable-doubt''}$. The jury's
context is $b = \text{``prosecution-narrative''}$. The attorney chooses signal
$s = \text{``alibi-evidence''}$. The language game payoff is
$\rs(\text{``alibi-evidence''} \compose \text{``prosecution-narrative''},
\text{``reasonable-doubt''})$. Crucially, the emergence term
$\emerge(s, b, a)$ captures how the evidence \emph{interacts} with the
prosecution's narrative to create reasonable doubt---an effect that neither
the evidence alone nor the narrative alone would produce.
\end{example}

\begin{example}[Making a Joke]
A comedian has intent $a = \text{``humor''}$ and audience context
$b = \text{``shared-assumptions''}$. The joke $s$ is structured as
setup + punchline. The setup primes the audience's context, and the
punchline produces emergence: $\emerge(s, b, a) \gg 0$. The humor
\emph{is} the emergence---the nonlinear interaction between the
punchline and the audience's primed expectations. A joke explained
($s = a$, direct transmission with zero emergence) is ``not funny''
precisely because the emergence term vanishes. This is why humor is
context-dependent: the same joke produces different emergence with
different audiences.
\end{example}

\begin{game}[The Translation Game]
A translator has a source text $a$ in one language (one ideatic space
$\IS_1$) and must produce a target signal $s$ in another language
(another ideatic space $\IS_2$) for a reader with context $b \in \IS_2$.
The translation ``succeeds'' when $\rs_2(s \compose b, T(a)) > 0$,
where $T : \IS_1 \to \IS_2$ is the semantic mapping between languages.

The impossibility of ``perfect translation'' is a corollary of the
emergence structure: even if $T$ preserves resonance ($\rs_1(a, c)
= \rs_2(T(a), T(c))$), the emergence in the target language
$\emerge_2(s, b, T(a))$ depends on the target reader's context $b$,
which has no counterpart in the source language. Every translation is a
new meaning game, not a mechanical transfer.
\end{game}


\section{Family Resemblance}

``Consider for example the proceedings that we call `games.' \ldots What is
common to them all? --- Don't say: `There must be something common, or they
would not be called ``games''\,' --- but \emph{look and see} whether there
is anything common to all.'' (\emph{PI} \S66)

Wittgenstein's concept of \emph{Familien\"ahnlichkeit} (family resemblance)
is his alternative to the classical theory of concepts-as-definitions. Instead
of a set of necessary and sufficient conditions, family resemblance posits a
network of overlapping similarities.

\begin{definition}[Family Resemblance]
\label{def:family-resemblance}
Ideas $a_1, \ldots, a_n \in \IS$ form a \emph{family} if there exists a
chain of positive pairwise resonances:
\[
\rs(a_i, a_{i+1}) > 0 \quad \text{for each } i = 1, \ldots, n-1.
\]
The \emph{family diameter} is the number of links in the longest
minimal chain.
\end{definition}

\begin{theorem}[Non-Transitivity of Resemblance]
\label{thm:non-transitive}
Family resemblance is \emph{not} transitive in general. There exist ideatic
spaces and ideas $a, b, c$ with $\rs(a, b) > 0$ and $\rs(b, c) > 0$ but
$\rs(a, c) \leq 0$.
\end{theorem}

\begin{proof}
We proceed by explicit construction. Consider the ideatic space
$\IS = \{e_1, e_2, e_3, \void\}$ with
\[
\rs(e_1, e_2) = 1, \quad \rs(e_2, e_3) = 1, \quad \rs(e_1, e_3) = -1.
\]
Then $e_1$ resembles $e_2$, $e_2$ resembles $e_3$, but $e_1$ and $e_3$
\emph{anti}-resonate. Board games resemble card games; card games resemble
drinking games; but board games may have nothing in common with---or even
oppose---drinking games.

The key mathematical observation is that the resonance function $\rs$
need not satisfy the triangle inequality. In a metric space,
$d(a,c) \leq d(a,b) + d(b,c)$ would force transitivity of closeness.
But $\rs$ is not a metric---it is a \emph{bilinear form} (as proved in
Volume~I), and bilinear forms can have indefinite signature. The non-transitivity
of resemblance is a direct consequence of the indefiniteness of $\rs$.
This is why the resonance function can capture what no distance function
can: the phenomenon of ideas that are similar-via-intermediary but
dissimilar directly.
\end{proof}

\begin{interpretation}[Contrast with Prototype Theory]
\emph{Classical concept theory} (Aristotle) defines categories by
necessary and sufficient conditions. \emph{Prototype theory} (Rosch, 1975)
replaces this with graded membership around a central exemplar.
\emph{Family resemblance} (Wittgenstein) rejects even the prototype: there
is no central member, only overlapping similarities.

In IDT, prototype theory corresponds to defining categories via resonance
with a fixed prototype $p$: idea $a$ belongs to the category iff
$\rs(a, p) > \theta$ for some threshold $\theta$. This is a \emph{linear}
classification---it uses only pairwise resonance.

Family resemblance, by contrast, requires the \emph{graph} structure of
the resonance function: the chain $\rs(a_i, a_{i+1}) > 0$ defines a
\emph{path} through the ideatic space, and the non-transitivity theorem
shows that path-connectedness does not imply direct resemblance. This
graph-theoretic perspective is richer than prototype theory and corresponds
precisely to Wittgenstein's ``complicated network of similarities
overlapping and criss-crossing'' (\emph{PI} \S66).
\end{interpretation}

\begin{definition}[Resonance Profile]
\label{def:resonance-profile}
The \emph{resonance profile} of $a \in \IS$ is the function
\[
\rho_a : \IS \to \R, \quad \rho_a(c) := \rs(a, c).
\]
Two ideas have the same meaning (in the Wittgensteinian sense) if and only
if they have the same resonance profile: $a \sim b$ iff $\rho_a = \rho_b$.
\end{definition}

\begin{theorem}[Meaning Is Use]
\label{thm:meaning-is-use}
The equivalence relation $a \sim b \iff \rho_a = \rho_b$ satisfies:
\begin{enumerate}
\item $a \sim a$ (reflexivity),
\item $a \sim b \implies b \sim a$ (symmetry),
\item $a \sim b \land b \sim c \implies a \sim c$ (transitivity),
\item $\void \not\sim a$ for any $a \neq \void$ with $w(a) > 0$ (silence has unique meaning).
\end{enumerate}
\end{theorem}

\begin{proof}
(1)--(3) are immediate from equality of functions. For (4), if $a \neq \void$
then $\rho_a(a) = \rs(a, a) = w(a) > 0$, but $\rho_\void(a) = \rs(\void, a) = 0$
by axiom R2. So $\rho_a \neq \rho_\void$.

The deeper point is that the equivalence relation $\sim$ is \emph{not}
identity: two syntactically distinct ideas can have identical meaning.
This is the formal correlate of synonymy. Moreover, the quotient
$\IS / {\sim}$ is the space of \emph{meanings}---the abstract semantic
space obtained by identifying all ideas with the same use.
The resonance function descends to this quotient:
$\rs([a], [b]) := \rs(a, b)$ is well-defined (one checks that
$a \sim a'$ and $b \sim b'$ imply $\rs(a, b) = \rs(a', b')$).
This quotient construction gives precise mathematical content to
Wittgenstein's claim that ``the meaning of a word is its use in the
language''---meaning is literally the equivalence class under the
resonance-profile relation.
\end{proof}

\begin{interpretation}[Habermas's Communicative Rationality]
J\"urgen Habermas's theory of \emph{communicative action} (1981) holds
that rational discourse presupposes three validity claims: truth
(propositional content), rightness (normative appropriateness), and
sincerity (authentic self-expression). A discourse is \emph{rational}
when these claims can be redeemed through argument.

In IDT, Habermas's validity claims map to dimensions of the resonance
profile:
\begin{itemize}
\item \textbf{Truth:} $\rs(a, \text{facts}) > 0$---the utterance resonates
with the factual context.
\item \textbf{Rightness:} $\rs(a, \text{norms}) > 0$---the utterance
resonates with the normative context.
\item \textbf{Sincerity:} $\rs(a, a_{\text{intent}}) > 0$---the utterance
resonates with the speaker's actual intent (non-deception).
\end{itemize}
Habermas's ``ideal speech situation''---a discourse free from domination
---corresponds to a meaning game where the optimal signal is the honest
signal $s = a$ (Corollary~\ref{cor:honest-signal}): no strategic
distortion is advantageous because the context is fully cooperative
($\emerge(a, b, a) \geq 0$ for all $a, b$). The pathology of
``systematically distorted communication'' (ideology, propaganda) occurs
when emergence is engineered to be negative for truthful signals but
positive for deceptive ones.
\end{interpretation}

\begin{lstlisting}[caption={Lean 4: Meaning-is-use equivalence}]
/-- Two ideas have the same meaning iff identical resonance profiles -/
def sameMeaning {I : Type*} [S : IdeaticSpace I] (a b : I) : Prop :=
  ∀ c : I, rs a c = rs b c

theorem sameMeaning_refl {I : Type*} [S : IdeaticSpace I] (a : I) :
    sameMeaning a a :=
  fun _ => rfl

theorem sameMeaning_symm {I : Type*} [S : IdeaticSpace I]
    {a b : I} (h : sameMeaning a b) :
    sameMeaning b a :=
  fun c => (h c).symm

theorem sameMeaning_trans {I : Type*} [S : IdeaticSpace I]
    {a b c : I} (h1 : sameMeaning a b) (h2 : sameMeaning b c) :
    sameMeaning a c :=
  fun d => (h1 d).trans (h2 d)

theorem void_unique_meaning {I : Type*} [S : IdeaticSpace I]
    {a : I} (ha : a ≠ void) : ¬sameMeaning a void := by
  intro h
  have := h a
  rw [rs_void_left'] at this
  exact absurd (rs_nondegen' a (le_antisymm (le_of_eq this.symm)
    (rs_self_nonneg' a))) ha
\end{lstlisting}


\section{Forms of Life}

``If a lion could speak, we could not understand him.'' (\emph{PI}, p.~223)

This is among Wittgenstein's most quoted and most debated remarks. In IDT, it
admits a precise formulation.

\begin{definition}[Form of Life]
\label{def:form-of-life}
A \emph{form of life} is an element $b \in \IS$ that serves as a shared
context for a community of speakers. Two speakers share a form of life if
their contexts $b_1, b_2$ satisfy $\rs(b_1, b_2) > 0$.
\end{definition}

\begin{theorem}[The Lion Theorem]
\label{thm:lion}
If the sender's form of life $b_S$ and the receiver's form of life $b_R$
satisfy $\rs(b_S, b_R) = 0$ (orthogonal forms of life), then for any signal
$s$ with $\rs(s, a) > 0$ (a signal that directly resonates with the intent),
the sender's payoff satisfies:
\[
u_S(s) = \rs(s, a) + \emerge(s, b_R, a).
\]
The preexisting alignment $\rs(b_R, a) = 0$ contributes nothing. Communication
depends entirely on direct resonance and emergence. If emergence is also zero
(the Augustinian case), then $u_S(s) = \rs(s, a)$: communication reduces
to transmission without understanding.
\end{theorem}

\begin{proof}
By the sender payoff decomposition (Theorem~\ref{thm:sender-decomp}):
$u_S(s) = \rs(s, a) + \rs(b_R, a) + \emerge(s, b_R, a)$.
Orthogonality of forms of life means $\rs(b_R, a) = 0$ (the receiver's form
of life has no resonance with the sender's intent, which is embedded in
$b_S$). So $u_S(s) = \rs(s, a) + \emerge(s, b_R, a)$.
\end{proof}

\begin{interpretation}
The lion can speak; we hear the sounds; but we cannot understand because our
form of life has zero resonance with the lion's intent. Even if the lion
says exactly what it means ($s = a$, so $\rs(s, a) = w(a) > 0$), the
received idea $s \compose b_R$ resonates with $a$ only through the direct
hit and emergence---the form of life contributes nothing. Understanding
requires shared context, and shared context means positive resonance between
forms of life.
\end{interpretation}

\begin{example}[Academic Peer Review]
Reviewer $R$ has form of life $b_R = \text{``analytic-philosophy''}$. Author
$A$ submits a paper with intent $a$ embedded in form of life
$b_A = \text{``continental-philosophy''}$. If $\rs(b_A, b_R) \approx 0$, the
reviewer ``cannot understand the lion'': the paper's meaning, composed with
the reviewer's context, produces no resonance with the author's intent. The
emergence term $\emerge(s, b_R, a)$ represents \emph{productive
misunderstanding}---the reviewer may find meanings the author never intended,
but these accidental resonances are not understanding.
\end{example}

\begin{example}[Religious Rhetoric]
A preacher speaks to a congregation sharing form of life $b_\text{faith}$.
The same sermon delivered to an atheist audience (form of life
$b_\text{secular}$) produces entirely different emergence. The words are
identical ($s$ is the same), but $\emerge(s, b_\text{faith}, a) \neq
\emerge(s, b_\text{secular}, a)$. The ``same'' sermon means different things
in different forms of life---this is the precise sense in which meaning is
use, not reference.
\end{example}

\begin{remark}[Incommensurability and Kuhn's Paradigm Shifts]
Thomas Kuhn's \emph{The Structure of Scientific Revolutions} (1962)
argues that scientists working in different paradigms are, in a sense,
working in ``different worlds.'' The concept of paradigm
incommensurability maps directly onto the Lion Theorem: two paradigms
$b_1$ (Newtonian mechanics) and $b_2$ (relativistic mechanics) may have
$\rs(b_1, b_2) \approx 0$, making cross-paradigm communication nearly
impossible. A Newtonian hearing a relativist speak about ``mass''
hears the word but composes it with a different context; the emergence
is different; the meaning is different. The scientific revolution is
not just a change in \emph{beliefs} but a change in \emph{form of
life}, which changes the emergence landscape of all subsequent
communication.

Kuhn's critics (Popper, Lakatos, Laudan) argued that incommensurability
is too strong: scientists \emph{can} communicate across paradigms,
even if imperfectly. IDT accommodates this nuance: if
$\rs(b_1, b_2) > 0$ (even slightly), communication is possible but
costly. The Lion Theorem's $\rs(b_S, b_R) = 0$ is the extreme case.
In practice, forms of life are rarely perfectly orthogonal; there is
always some residual resonance, some shared humanity, that enables
the conversation to continue---however imperfectly.
\end{remark}


\section{Rule-Following and Private Language}

``This was our paradox: no course of action could be determined by a rule,
because every course of action can be made out to accord with the rule.''
(\emph{PI} \S201)

The rule-following paradox is Wittgenstein's deepest challenge to the idea of
determinate meaning. If meaning is fixed by rules, and rules require
interpretation, and interpretations require further rules \ldots then meaning
seems to dissolve in infinite regress.

IDT dissolves this paradox by \emph{not grounding meaning in rules at all}.
Meaning is resonance, and resonance is a primitive: $\rs(a, b)$ is a
real number, not an interpretation of a rule. There is no regress because
there is no rule to interpret---there is only the resonance function, which is
given as part of the structure.

\begin{theorem}[No Private Language]
\label{thm:no-private-language}
In any ideatic space $(\IS, \compose, \void, \rs)$, if idea $a$ is
``private'' in the sense that $\rs(a, c) = 0$ for all $c \neq a$, then
the language game $(\IS, \compose, \void, \rs, a, b)$ has value
$u_S(s) = \rs(s, a) + \emerge(s, b, a)$ for all $s$---and the only signal
that can succeed is one that directly resonates with $a$ or produces emergence
with the receiver's context. But if $a$ is truly private
($\rs(s, a) = 0$ for all $s \neq a$), then the only effective signal is
$s = a$ itself: private meaning can only be communicated by direct
transmission, not by indirect reference.
\end{theorem}

\begin{proof}
If $\rs(s, a) = 0$ for all $s \neq a$, then $u_S(s) = \emerge(s, b, a)$
for $s \neq a$. By the emergence bound
$|\emerge(s, b, a)|^2 \leq w(s \compose b) \cdot w(a)$, emergence is
bounded. But the only signal with \emph{guaranteed} positive payoff is
$s = a$, which gives $u_S(a) = w(a) + \emerge(a, b, a) \geq w(a) - \sqrt{w(a \compose b) \cdot w(a)}$. A truly private idea can only be communicated
by saying exactly itself---there are no ``indirect routes'' through the
ideatic space.
\end{proof}

\begin{interpretation}
Wittgenstein's private language argument (\emph{PI} \S\S243--315) concludes
that a purely private language is impossible. In IDT, the corresponding
result is that a purely private idea (one with zero resonance with everything
else) can only be communicated by direct transmission. It cannot participate
in the rich game of indirect reference, metaphor, and context-dependent
meaning that constitutes natural language. A beetle in a box that no one
can see (\emph{PI} \S293) is an idea with zero resonance: it ``drops out of
consideration as irrelevant.''
\end{interpretation}

\begin{remark}[The Beetle in the Box, Formalized]
Wittgenstein's beetle-in-the-box thought experiment asks us to imagine
that everyone has a box with something in it called a ``beetle,'' but
no one can look in anyone else's box. In IDT, the beetle is an idea
$p$ with $\rs(p, c) = 0$ for all $c \neq p$. As proved in Volume~I
(the transparency theorem), when an idea is transparent to all observers,
it contributes nothing to emergence:
$\emerge(c, p, d) = 0$ for all $c, d$ when $p$ is private. The beetle
``drops out of consideration'' because it generates zero emergence---it
adds no nonlinear meaning to any composition. The box might as well
be empty.
\end{remark}

\begin{remark}[Kripke's Skeptical Solution]
Saul Kripke's \emph{Wittgenstein on Rules and Private Language} (1982)
reads the rule-following argument as a \emph{skeptical paradox}: there is
no fact about what rule a speaker is following. Kripke's ``skeptical
solution'' appeals to community agreement: meaning is constituted by the
community's dispositions to accept or reject applications.

In IDT, Kripke's community-based account corresponds to defining meaning
through \emph{aggregate resonance profiles}. The meaning of $a$ is not
$\rho_a(c)$ for any single observer $c$, but the profile
$\rho_a^{\mathcal{C}} := \sum_{c \in \mathcal{C}} \rho_a(c)$ across a
community $\mathcal{C}$. This aggregate profile is determinate (no
skeptical paradox) because the resonance function is given as part of the
structure---there is nothing to ``interpret.'' IDT resolves Kripke's
paradox not by finding a meaning-constituting fact, but by replacing facts
with resonances: meaning is not a property an idea \emph{has}, but a
relation it \emph{stands in} with the community.
\end{remark}


\section{Agreement in Judgments}

``It is what human beings say that is true and false; and they agree in the
\emph{language} they use. That is not agreement in opinions but in form of
life.'' (\emph{PI} \S241)

\begin{definition}[Agreement in Form of Life]
\label{def:agreement-form}
Two agents with contexts $b_1, b_2 \in \IS$ \emph{agree in form of life}
if $\rs(b_1, b_2) > 0$. They \emph{agree in judgments about} $a$ if
$\rs(b_1, a) \cdot \rs(b_2, a) > 0$ (both resonate positively or both
resonate negatively with $a$).
\end{definition}

\begin{theorem}[Agreement Decomposition]
\label{thm:agreement-decomp}
The correlation of two receivers' responses to a shared signal $s$ is:
\[
\rs(s \compose b_1, a) \cdot \rs(s \compose b_2, a) =
[\rs(s, a) + \rs(b_1, a) + \emerge(s, b_1, a)] \cdot
[\rs(s, a) + \rs(b_2, a) + \emerge(s, b_2, a)].
\]
Agreement (both terms positive) requires that the sum of direct resonance,
preexisting alignment, and emergence be positive for both receivers.
\end{theorem}

\begin{proof}
Direct application of the sender payoff decomposition to each receiver.
\end{proof}

\begin{interpretation}
Two people can agree on the meaning of a sentence even though they arrive
at this agreement by different routes: one through preexisting alignment
(ethos), the other through emergence (pathos). What matters is that the
\emph{total} effect is positive for both. This is Wittgenstein's ``agreement
in form of life'' made precise: not agreement on any particular content, but
agreement in the \emph{pattern of resonance}.
\end{interpretation}

\begin{remark}[Disagreement as Negative Emergence]
When two receivers disagree about a shared signal---one finding it meaningful,
the other not---the disagreement has a precise source. Writing out the product:
\[
\rs(s \compose b_1, a) \cdot \rs(s \compose b_2, a) < 0
\]
means one factor is positive and the other negative. This can arise from
different preexisting alignments ($\rs(b_1, a)$ vs.\ $\rs(b_2, a)$) or
different emergence responses ($\emerge(s, b_1, a)$ vs.\ $\emerge(s, b_2, a)$).
The latter is the more interesting case: the \emph{same signal}, in
\emph{different contexts}, creates opposite meanings. This is why the same
joke can be funny to one audience and offensive to another, why the same
political slogan can inspire one group and alienate another. The
signal hasn't changed; the context has.

In hermeneutics, this is Gadamer's insight that understanding is always
``prejudiced'' by one's horizon (Gadamer 1960). IDT makes the prejudice
precise: it is the receiver's context $b_i$, which enters both the
direct resonance term $\rs(b_i, a)$ and the emergence term
$\emerge(s, b_i, a)$. The ``fusion of horizons'' that Gadamer describes
is the condition under which two receivers with different contexts
nevertheless achieve positive correlation: they agree despite their
differences because the signal's emergence is robust to contextual
variation.
\end{remark}


\section{Game-Theoretic Contrast: Why Not Classical Signaling Games?}

The classical signaling game (Spence 1973, Crawford--Sobel 1982) models
communication as follows: a sender of type $\theta$ sends signal $s$; a
receiver takes action $a(s)$; payoffs $u_S(\theta, s, a(s))$ and
$u_R(\theta, s, a(s))$ depend on type, signal, and action.

\begin{remark}[What IDT Adds to Classical Signaling]
The IDT meaning game differs from the classical signaling game in three
fundamental ways:
\begin{enumerate}
\item \textbf{Composition, not action.} The receiver does not choose an
``action''---the receiver \emph{composes} the signal with their context. The
``action'' is automatic: it is the act of understanding itself.
\item \textbf{Emergence.} The payoff is not a fixed function of signal and
type. It depends on the \emph{emergence} $\emerge(s, b, a)$, which captures
the genuinely novel meaning created by the interaction of signal and context.
Classical games have no analogue of emergence.
\item \textbf{Resonance, not utility.} The payoff function $\rs$ is not a
preference ordering over outcomes. It is a \emph{resonance}: a measure of
how well ideas fit together. This makes meaning games inherently
non-zero-sum in a way that classical games are not.
\end{enumerate}
\end{remark}

Von Neumann and Morgenstern's \emph{Theory of Games and Economic Behavior}
(1944) assumed that players have fixed utility functions over outcomes. Nash
(1950) proved the existence of equilibria under this assumption. But meaning
games violate the assumption: the ``outcome'' of composing $s$ with $b$ is
not fixed by $s$ and $b$ separately---it depends on the emergence, which is
a property of the composition itself. This is why IDT needs its own
equilibrium theory, developed in Chapter~\ref{ch:equilibrium}.

\begin{interpretation}[Three Traditions, One Framework]
Three intellectual traditions converge in this chapter:

\textbf{(a) Classical game theory} (von Neumann, Nash, Selten) models
strategic interaction with fixed payoff structures. It assumes
\emph{separability}: outcomes decompose into individual contributions.
Emergence violates separability.

\textbf{(b) Analytic philosophy of language} (Wittgenstein, Austin, Grice,
Lewis, Brandom) analyzes meaning through use, speech acts, conversational
implicature, convention, and inferential roles. Each captures one facet of
the resonance function: Austin's speech acts are strategies in meaning
games; Grice's maxims are emergence constraints (Chapter~\ref{ch:deception});
Lewis's conventions are Nash equilibria with $\emerge = 0$; Brandom's
inferential roles are resonance profiles.

\textbf{(c) Continental philosophy of dialogue} (Buber, Gadamer, Habermas,
Bakhtin) emphasizes the irreducibly relational character of meaning. Buber's
I-Thou relation, Gadamer's ``fusion of horizons,'' Habermas's communicative
action, Bakhtin's dialogism---all describe situations where the whole
exceeds the sum of parts. In IDT, this excess is precisely the emergence
term.

What each tradition \emph{lacks}, the others supply. Game theory lacks a
theory of meaning creation. Philosophy of language lacks a theory of
strategic interaction. Continental philosophy lacks mathematical precision.
IDT provides all three simultaneously: strategic interaction
(\S\ref{ch:equilibrium}--\ref{ch:two-sided}), meaning creation (through
$\emerge$), and mathematical precision (through the ideatic space axioms).
\end{interpretation}



% ================================================================
% CHAPTER 2: THE BASIC MEANING GAME
% ================================================================
\chapter{The Basic Meaning Game}
\label{ch:basic-game}

\epigraph{``The theory of games of strategy is a mathematical theory that
deals with the rational behavior of participants in situations of conflict
and cooperation.''}
{---John von Neumann \& Oskar Morgenstern, \emph{Theory of Games and
Economic Behavior}, 1944}


\section{Players, Strategies, Payoffs}

Von Neumann's formalization of games begins with three primitives: the set of
players, the set of strategies available to each, and a payoff function.
We now construct the meaning game from the primitives of IDT.

\begin{definition}[The Basic Meaning Game]
\label{def:basic-meaning-game}
A \emph{basic meaning game} $\Gamma(a, b)$ consists of:
\begin{itemize}
\item \textbf{Players:} A sender $S$ and a receiver $R$.
\item \textbf{Sender's type:} An idea $a \in \IS$ (intent).
\item \textbf{Receiver's context:} An idea $b \in \IS$ (form of life).
\item \textbf{Sender's strategy set:} $\IS$ (the sender may say anything).
\item \textbf{Receiver's strategy:} Passive composition: received idea $= s \compose b$.
\item \textbf{Payoffs:}
\begin{align}
u_S(s) &:= \rs(s \compose b, a), \label{eq:sender-payoff} \\
u_R(s) &:= w(s \compose b) - w(s) = \rs(s \compose b, s \compose b) - \rs(s, s). \label{eq:receiver-payoff}
\end{align}
\end{itemize}
\end{definition}

\begin{remark}[Comparison with Von Neumann--Morgenstern]
In a classical game, each player's strategy set is an abstract set, and
the payoff function is a map from strategy profiles to $\R$. In the meaning
game:
\begin{itemize}
\item The strategy set is the ideatic space $\IS$ itself---strategies
\emph{are} ideas.
\item The payoff function is derived from the resonance function---payoffs
\emph{are} resonances.
\item The receiver has no strategic choice: composition is automatic.
\end{itemize}
This makes the basic meaning game a \emph{one-player optimization problem}
embedded in a two-player structure. The game-theoretic richness emerges when
we extend to two-sided games (Chapter~\ref{ch:two-sided}) and repeated games.
\end{remark}

\begin{lstlisting}[caption={Lean 4: Meaning game payoffs in IDT v3}]
variable {I : Type*} [S : IdeaticSpace I]
open IdeaticSpace

/-- Sender payoff: resonance of received idea with intent -/
noncomputable def senderPayoff (s a b : I) : ℝ :=
  rs (compose s b) a

/-- Receiver payoff: enrichment from receiving signal -/
noncomputable def receiverPayoff (s b : I) : ℝ :=
  rs (compose s b) (compose s b) - rs s s

/-- Receiver payoff is non-negative by E4 -/
theorem receiverPayoff_nonneg (s b : I) :
    receiverPayoff s b ≥ 0 := by
  unfold receiverPayoff
  linarith [S.compose_enriches s b]
\end{lstlisting}

\begin{example}[Political Debate]
In a political debate, candidate $S$ has intent $a =
\text{``lower-taxes-good''}$. The audience's context is $b =
\text{``current-economic-anxiety''}$. The candidate's signal $s =
\text{``I-will-cut-your-taxes''}$ produces payoff
$u_S(s) = \rs(s \compose b, a)$. The receiver payoff $u_R(s) = w(s \compose b) - w(s)$ measures how much the audience's cognitive state is enriched
by hearing the promise \emph{in the context of their economic anxiety}.
By axiom E4, $u_R \geq 0$: the audience is always enriched, even by a lie.
\end{example}

\begin{example}[Advertising]
A company has intent $a = \text{``buy-our-product''}$. The consumer's context
is $b = \text{``daily-life-needs''}$. The ad signal $s =
\text{``life-is-better-with-X''}$ succeeds when $\rs(s \compose b, a) > 0$:
the composition of the ad with the consumer's daily-life context resonates
with the purchase intent. A \emph{clever} ad exploits emergence:
$\emerge(s, b, a)$ is large because the ad interacts with the consumer's
existing needs to create resonance with purchase that neither the ad nor
the needs would produce alone.
\end{example}


\section{The Fundamental Decomposition}

\begin{theorem}[Sender Payoff Decomposition]
\label{thm:sender-decomp}
For any signal $s$ in game $\Gamma(a, b)$:
\[
u_S(s) = \underbrace{\rs(s, a)}_{\text{direct resonance}} +
\underbrace{\rs(b, a)}_{\text{preexisting alignment}} +
\underbrace{\emerge(s, b, a)}_{\text{emergence}}.
\]
\end{theorem}

\begin{proof}
By definition of emergence:
$\emerge(s, b, a) = \rs(s \compose b, a) - \rs(s, a) - \rs(b, a)$.
Rearranging: $\rs(s \compose b, a) = \rs(s, a) + \rs(b, a) + \emerge(s, b, a)$.

The mathematical content is more than algebraic rearrangement. The
emergence decomposition says that the resonance of a composite idea
$s \compose b$ with any observer $a$ has exactly three sources. The first,
$\rs(s, a)$, depends only on the signal and the observer---it is the
\emph{context-free} contribution. The second, $\rs(b, a)$, depends only on
the receiver's context and the observer---it is the \emph{signal-free}
contribution. The third, $\emerge(s, b, a)$, depends on all three
simultaneously---it is the irreducibly \emph{triadic} component that
cannot be attributed to any pair of ideas alone.

This triadic structure is the key mathematical difference between IDT and
all preceding frameworks. Shannon's information theory is dyadic (sender-receiver).
Classical game theory is dyadic (strategy-payoff). Even Peirce's semiotics is
triadic (sign-object-interpretant) but lacks the quantitative emergence term.
The sender payoff decomposition makes the triadic structure \emph{computable}.
\end{proof}

\begin{lstlisting}[caption={Lean 4: Sender payoff decomposition}]
theorem senderPayoff_decomp (s a b : I) :
    senderPayoff s a b =
    rs s a + rs b a + emergence s b a := by
  unfold senderPayoff emergence; ring
\end{lstlisting}

This decomposition is the central equation of the book. It says that the
sender's payoff has exactly three sources, which we now analyze.

\begin{definition}[The Three Components of Communication]
\label{def:three-components}
In the sender payoff decomposition $u_S(s) = \rs(s,a) + \rs(b,a) + \emerge(s,b,a)$:
\begin{enumerate}
\item \textbf{Exchange} $:= \rs(s, a) + \rs(b, a)$. This is the
\emph{linear} part: what would happen if composition were additive
($\emerge \equiv 0$). It consists of the signal's direct resonance with the
intent, plus the receiver's preexisting alignment.
\item \textbf{Emergence} $:= \emerge(s, b, a)$. This is the \emph{nonlinear}
part: the genuinely new resonance created by the composition. It is the
surplus (or deficit) that arises because meaning is not additive.
\end{enumerate}
\end{definition}

\begin{theorem}[Exchange + Emergence Decomposition]
\label{thm:exchange-emergence}
Every sender payoff decomposes as:
\[
u_S(s) = \text{Exchange}(s, a, b) + \text{Emergence}(s, b, a),
\]
where the exchange is predictable from pairwise resonances and the emergence
is the unpredictable surplus.
\end{theorem}

\begin{proof}
This is a restatement of Theorem~\ref{thm:sender-decomp}.
\end{proof}

\begin{interpretation}[The Dual Nature of Communication]
The exchange + emergence decomposition reveals that every act of communication
has two aspects:
\begin{enumerate}
\item \textbf{Information transfer} (exchange): the signal carries resonance
from sender to receiver. This is the Shannon channel---predictable,
analyzable, decomposable.
\item \textbf{Meaning creation} (emergence): the signal interacts with the
receiver's context to create something new. This is what Shannon's theory
misses---unpredictable, irreducible, genuinely creative.
\end{enumerate}
Shannon (1948) modeled communication as exchange alone. Wittgenstein insisted
that meaning is more than information transfer. IDT quantifies the ``more'':
it is the emergence term $\emerge(s, b, a)$.
\end{interpretation}

\begin{remark}[The Shannon--Wittgenstein Divide]
The history of communication theory since 1948 can be read as a struggle
between Shannon's quantitative but meaning-free framework and
Wittgenstein's meaning-rich but non-quantitative philosophy. Shannon
himself warned that his theory ``has nothing to do with meaning''---a
remark that has been both celebrated (by engineers) and lamented (by
humanists).

IDT resolves this divide by \emph{containing both}. The exchange term
$\rs(s, a) + \rs(b, a)$ is Shannon's information channel: it captures
the predictable, decomposable aspect of communication. The emergence term
$\emerge(s, b, a)$ is Wittgenstein's ``more'': the meaning that arises
from use, context, and form of life. The sender payoff decomposition
is the equation that unifies the two traditions:
\[
\underbrace{u_S(s)}_{\text{total meaning}} =
\underbrace{\rs(s,a) + \rs(b,a)}_{\text{Shannon}} +
\underbrace{\emerge(s,b,a)}_{\text{Wittgenstein}}.
\]
Every prior theory of communication is a special case of this equation.
\end{remark}


\section{The Honest Signal and Strategic Alternatives}

\begin{corollary}[The Honest Signal]
\label{cor:honest-signal}
If the sender uses signal $s = a$ (says exactly what they mean):
\[
u_S(a) = w(a) + \rs(b, a) + \emerge(a, b, a).
\]
The honest signal earns at least $w(a)$ (the weight of the idea itself)
plus preexisting alignment and emergence.
\end{corollary}

\begin{theorem}[Optimal Signal Characterization]
\label{thm:optimal-signal}
A signal $s^*$ is optimal iff it maximizes $\rs(s, a) + \emerge(s, b, a)$
over $s \in \IS$. The preexisting alignment $\rs(b, a)$ is a constant and
does not affect the choice.
\end{theorem}

\begin{proof}
$u_S(s) = \rs(s,a) + \rs(b,a) + \emerge(s,b,a)$. Since $\rs(b,a)$ does not
depend on $s$, maximizing $u_S(s)$ is equivalent to maximizing
$\rs(s,a) + \emerge(s,b,a)$.
\end{proof}

\begin{remark}[Strategic Dishonesty]
The optimal signal need not be the honest signal $s = a$. If there exists
$s^* \neq a$ with $\rs(s^*, a) + \emerge(s^*, b, a) > w(a) + \emerge(a, b, a)$,
then $s^*$ outperforms honesty. This occurs when the emergence from
a strategically chosen signal exceeds the weight advantage of honesty.
Rhetoric, persuasion, and propaganda all exploit this possibility.
\end{remark}

\begin{example}[Therapy]
A therapist has intent $a = \text{``patient-should-recognize-pattern''}$.
The patient's context is $b = \text{``patient's-defenses''}$. Direct
statement $s = a$ may fail: $\emerge(a, b, a) < 0$ because the patient's
defenses interact destructively with direct confrontation. Instead, the
therapist chooses an indirect signal $s = \text{``tell-me-about-your-week''}$
with $\rs(s, a)$ small but $\emerge(s, b, a)$ large: the question, composed
with the patient's defenses, creates emergent resonance with the therapeutic
intent.
\end{example}


\section{The Cocycle Constraint on Emergence}

The emergence function is not arbitrary: it satisfies a deep algebraic
constraint inherited from the associativity of composition.

\begin{theorem}[Cocycle Condition]
\label{thm:cocycle}
For all $a, b, c, d \in \IS$:
\[
\emerge(a, b \compose c, d) + \emerge(b, c, d)
= \emerge(a \compose b, c, d) + \emerge(a, b, d).
\]
\end{theorem}

\begin{proof}
Expand both sides using the definition
$\emerge(x, y, d) = \rs(x \compose y, d) - \rs(x, d) - \rs(y, d)$:

LHS $= [\rs(a \compose (b \compose c), d) - \rs(a, d) - \rs(b \compose c, d)]
+ [\rs(b \compose c, d) - \rs(b, d) - \rs(c, d)]
= \rs(a \compose b \compose c, d) - \rs(a, d) - \rs(b, d) - \rs(c, d)$.

RHS $= [\rs((a \compose b) \compose c, d) - \rs(a \compose b, d) - \rs(c, d)]
+ [\rs(a \compose b, d) - \rs(a, d) - \rs(b, d)]
= \rs(a \compose b \compose c, d) - \rs(a, d) - \rs(b, d) - \rs(c, d)$.

Both sides equal the same expression by associativity
$a \compose (b \compose c) = (a \compose b) \compose c$.

The deeper insight is that the cocycle condition is a \emph{consequence of
associativity}, not an independent axiom. The emergence function, though
nonlinear and apparently unconstrained, is in fact tightly governed by the
algebraic structure of composition. This is analogous to group cohomology:
in the cohomology of groups, a 2-cocycle $\omega(g, h)$ satisfies
$\omega(g, hk) + \omega(h, k) = \omega(gh, k) + \omega(g, h)$---exactly
the same algebraic form as our emergence cocycle. The emergence function
$\emerge$ is a 2-cocycle on the monoid $(\IS, \compose, \void)$ with
values in the abelian group $(\R, +)$, where the ``coefficients'' are
determined by the observer $d$.

This cohomological perspective has a profound consequence: the emergence
function is classified (up to coboundaries) by the second cohomology
group $H^2(\IS, \R)$. Two ideatic spaces with the same underlying monoid
but different resonance functions will have the same ``type'' of emergence
if and only if their emergence cocycles are cohomologous. The cocycle
condition is not just a constraint---it is a \emph{classification tool}.
\end{proof}

\begin{lstlisting}[caption={Lean 4: The cocycle condition}]
theorem cocycle_condition (a b c d : I) :
    emergence a (compose b c) d + emergence b c d =
    emergence (compose a b) c d + emergence a b d := by
  unfold emergence
  rw [compose_assoc']
  ring
\end{lstlisting}

\begin{interpretation}
The cocycle condition is the ``conservation law'' of emergence. It says that
emergence is \emph{path-independent} in a precise algebraic sense: the total
emergence from composing $a$, $b$, $c$ and observing against $d$ does not
depend on whether we first compose $a$ with $b$ and then with $c$, or first
compose $b$ with $c$ and then with $a$. This is a deep constraint: it means
that emergence, though nonlinear, is not \emph{arbitrary}. It has structure.
\end{interpretation}


\section{Receiver Payoff and Enrichment}

\begin{theorem}[Receiver Is Always Enriched]
\label{thm:receiver-enriched}
For all $s, b \in \IS$: $u_R(s) = w(s \compose b) - w(s) \geq 0$.
\end{theorem}

\begin{proof}
By axiom E4: $w(s \compose b) \geq w(s)$ (composition enriches).
Actually, E4 states $w(x \compose y) \geq w(x)$, so
$w(s \compose b) \geq w(s)$.
\end{proof}

\begin{remark}[Contrast with Classical Games]
In a classical game, a player can \emph{lose} from an interaction. In the
basic meaning game, the receiver \emph{always gains} (in weight). This is
the fundamental asymmetry of meaning games: receiving a signal never
decreases your weight. You can be deceived, manipulated, or confused, but
you are always enriched. This is why propaganda works: even a false signal
enriches the receiver, making the false idea ``heavier'' in their cognitive
state.
\end{remark}

\begin{interpretation}[The Asymmetry of Sender and Receiver]
The receiver-always-enriched theorem reveals a deep asymmetry between
the two roles in communication:

\textbf{(a) The sender bears all risk.} The sender's payoff
$u_S(s) = \rs(s \compose b, a)$ can be positive, zero, or negative.
Sending a signal that anti-resonates with the intent (due to negative
emergence) makes the sender \emph{worse off}. The sender gambles on
the emergence term.

\textbf{(b) The receiver is always safe.} The receiver's payoff
$u_R(s) = w(s \compose b) - w(s) \geq 0$ is guaranteed non-negative.
No matter how bad the signal, the receiver is at least as well off as
before. This is because weight (self-resonance) can only increase under
composition.

\textbf{(c) The new insight} is that this asymmetry explains the
fundamental power imbalance in communication. The receiver holds all the
cards: they can listen to any signal at no risk, while the sender must
carefully choose signals to avoid negative emergence. In negotiations,
this is the ``listener's advantage''---the party that listens more and
speaks less bears less risk. In journalism, this is why sources are
vulnerable and journalists are powerful: the source (sender) bears the
risk of negative emergence, while the journalist (receiver) is always
enriched by new information.

This also explains why \emph{attention} is the scarce resource in the
modern information economy. Receiving a signal is always beneficial
(the receiver is enriched), so the bottleneck is not willingness to
listen but \emph{capacity} to listen. In a world of unlimited signals,
the receiver's attention---not the sender's content---is the binding
constraint.
\end{interpretation}


\section{Costly Signaling and the Economics of Honest Communication}

The fundamental problem of communication theory since Spence (1973) is:
\emph{how can signals be credible?} If cheap talk is costless, why believe
anything? IDT provides a natural cost structure: the \emph{weight} of the
signal is its cost.

\begin{definition}[Signal Cost]
\label{def:signal-cost}
The \emph{cost} of signal $s$ is its self-resonance (weight):
\[
C(s) := w(s) = \rs(s, s).
\]
More substantial signals---those with higher weight---are more costly to
produce. Void is the only free signal: $C(\void) = 0$.
\end{definition}

\begin{theorem}[Costly Net Payoff]
\label{thm:costly-net}
The sender's \emph{net payoff} in a costly signaling game is:
\[
u_S^{\text{net}}(s) := u_S(s) - C(s) = \rs(s \compose b, a) - w(s).
\]
A signal is \emph{worth sending} if and only if $u_S^{\text{net}}(s) > 0$:
the resonance it creates exceeds its intrinsic cost.
\end{theorem}

\begin{proof}
We proceed by observing that the sender payoff decomposition gives
$u_S(s) = \rs(s, a) + \rs(b, a) + \emerge(s, b, a)$, so the net payoff is
\[
u_S^{\text{net}}(s) = \rs(s, a) + \rs(b, a) + \emerge(s, b, a) - w(s).
\]
For this to be positive, we need $\rs(s, a) + \emerge(s, b, a) > w(s) - \rs(b, a)$.
Since $\rs(s, a) \leq \sqrt{w(s) \cdot w(a)}$ by Cauchy--Schwarz, costly
signals must generate sufficient emergence to overcome their weight. This is
the formal content of ``credibility through cost'': a heavy signal ($w(s)$ large)
can only be worth sending if it creates correspondingly large emergence. Low-type
senders who lack the capacity to generate this emergence cannot profitably
imitate high-type signals.
\end{proof}

\begin{lstlisting}[caption={Lean 4: Costly signaling}]
/-- Signal cost is its self-resonance (weight) -/
noncomputable def signalCost (s : I) : ℝ := rs s s

/-- Void is free -/
theorem signalCost_void : signalCost (void : I) = 0 := rs_void_void

/-- Non-void signals cost strictly positive -/
theorem signalCost_pos (s : I) (h : s ≠ void) :
    signalCost s > 0 :=
  rs_pos_of_ne_void s h

/-- A signal is "worth sending" iff resonance exceeds cost -/
theorem signal_worth_sending_iff (s r : I) :
    senderPayoff s r > signalCost s ↔
    senderPayoff s r - signalCost s > 0 := by
  constructor <;> intro h <;> linarith
\end{lstlisting}

\begin{interpretation}[Spence's Job Market Signaling]
In Spence's (1973) model, a job applicant's education level serves as a
costly signal of ability. The key insight is that education is more
costly for low-ability workers, creating a \emph{separating equilibrium}
where only high-ability workers invest in education.

In IDT, this translates directly: a worker's ``ability'' is their type
$a$, and education is a signal $s$ with cost $w(s)$. The composition
$s \compose b$ (education in the employer's context) creates emergence
$\emerge(s, b, a)$ that differs by type. For high-ability workers, the
emergence is large (education genuinely interacts with their knowledge
to create value); for low-ability workers, the emergence is small.
The Spence separation condition becomes:
\[
\rs(s, a_H) + \emerge(s, b, a_H) - w(s) > 0 >
\rs(s, a_L) + \emerge(s, b, a_L) - w(s).
\]
IDT adds what Spence's model lacks: a theory of \emph{why} education
is differentially costly---it is because the emergence from composing
education with the employer's context \emph{differs by type}. Cost is
not exogenous; it is an emergence property.
\end{interpretation}


\section{The Babbling Equilibrium and Communicative Failure}

The babbling equilibrium is the canonical case of communicative failure:
the receiver ignores every signal, and the sender has no incentive to
communicate.

\begin{definition}[Babbling Receiver]
\label{def:babbling}
A receiver $b$ is \emph{babbling} if they ignore all signals:
\[
\rs(b, s \compose b) = \rs(b, b) \quad \text{for all } s \in \IS.
\]
Equivalently, the receiver payoff is independent of the signal.
\end{definition}

\begin{theorem}[Babbling Equilibrium Properties]
\label{thm:babbling-properties}
Under babbling:
\begin{enumerate}
\item Communication surplus reduces to the sender's net payoff alone.
\item Void signal achieves net payoff $0$, so silence is always weakly optimal.
\item The misunderstanding gap equals the sender's self-resonance.
\item Total social welfare reduces to $\rs(s, s) + 2\rs(b, b)$.
\end{enumerate}
\end{theorem}

\begin{proof}
Under babbling, $\rs(b, s \compose b) = w(b)$ for all $s$, so the receiver
payoff $u_R(s) = w(s \compose b) - w(s)$ is constant across signals---the
receiver's contribution to surplus is signal-independent. The communication
surplus therefore reduces to the sender's net gain.

For void, $u_S^{\text{net}}(\void) = \rs(\void \compose b, a) - 0 = \rs(b, a)$,
which is the sender's baseline. Any signal $s$ with
$u_S^{\text{net}}(s) < \rs(b, a)$ is dominated by silence. Under babbling, since
the receiver payoff is constant, the void signal is always weakly optimal:
$u_S(\void) = \rs(b, a) \geq u_S(s)$ whenever $\rs(s, a) + \emerge(s, b, a) \leq 0$.
\end{proof}

\begin{lstlisting}[caption={Lean 4: Babbling equilibrium}]
/-- A receiver is babbling if they ignore all signals -/
def isBabbling (r : I) : Prop :=
  ∀ s : I, receiverPayoff s r = receiverPayoff (void : I) r

/-- Under babbling, void is weakly optimal for sender -/
theorem babbling_void_optimal (r : I) (_h : isBabbling r) :
    senderNetPayoff (void : I) r = 0 :=
  senderNetPayoff_void_signal r

/-- Under babbling, communication surplus = sender's net -/
theorem babbling_surplus (r : I) (h : isBabbling r) (s : I) :
    communicationSurplus s r =
    senderNetPayoff s r + receiverPayoff (void : I) r := by
  unfold communicationSurplus
  rw [show receiverPayoff s r = receiverPayoff (void : I) r from h s]
\end{lstlisting}

\begin{interpretation}[Babbling as the Crawford--Sobel Limit]
Crawford and Sobel (1982) showed that in signaling games with
partially aligned preferences, the set of equilibria ranges from
fully revealing (many distinct messages) to \emph{babbling} (one
pooling message). The babbling equilibrium exists in \emph{every}
such game---it is the universal fallback when communication fails.

In IDT, the babbling condition $\rs(b, s \compose b) = w(b)$ for all $s$
gives this a precise geometric interpretation: the receiver's form of life
$b$ is \emph{orthogonal to all emergence}. No signal can create nonlinear
interaction with the receiver's context. The receiver is, in effect,
\emph{immune to meaning}---they hear signals as pure noise.

This connects to Wittgenstein's lion: the lion can speak, and we hear
sounds, but we cannot extract meaning because our form of life generates
no emergence with the lion's signals. The babbling equilibrium is the
\emph{game-theoretic name} for Wittgensteinian incomprehension.
\end{interpretation}

\begin{remark}[The Babbling-Silence Duality]
There is a revealing duality between babbling (receiver's failure) and
silence (sender's choice). When the receiver babbles, the sender's best
response is often silence: why speak if no one listens? When the sender
is silent, the receiver's best response is babbling: why listen if no one
speaks? This creates a \emph{communicative trap}: the void-void equilibrium
is self-reinforcing. Breaking out requires an exogenous shock---a signal
so powerful that it creates emergence even with a babbling receiver, or
a context shift that makes the receiver responsive.

As proved in the Lean formalization (T208), under babbling, the silence
gap and the babbling gap coincide: the sender's reason for silence is
exactly the receiver's reason for not listening.
\end{remark}



% ================================================================
% CHAPTER 3: EQUILIBRIUM IN MEANING GAMES
% ================================================================
\chapter{Equilibrium in Meaning Games}
\label{ch:equilibrium}

\epigraph{``An equilibrium point is an $n$-tuple such that each player's
mixed strategy maximizes his payoff if the strategies of the others are
held fixed.''}
{---John Nash, ``Non-Cooperative Games,'' 1951}


\section{Nash's Concept and Its Limits}

Nash's 1950 theorem guarantees that every finite game has a mixed-strategy
equilibrium. But meaning games are not finite: the strategy set is the entire
ideatic space $\IS$, which is infinite. Moreover, the payoff function depends
on the emergence term, which introduces nonlinearities that have no analogue
in classical games.

We therefore develop the equilibrium theory of meaning games from scratch,
using Nash's framework as a starting point but not as a constraint.

\begin{definition}[Best Response]
\label{def:best-response}
In the meaning game $\Gamma(a, b)$, signal $s^*$ is a \emph{best response}
for the sender if
\[
u_S(s^*) \geq u_S(s) \quad \text{for all } s \in \IS.
\]
Equivalently, $s^*$ maximizes $\rs(s, a) + \emerge(s, b, a)$ over $\IS$.
\end{definition}

\begin{definition}[Nash Equilibrium in Meaning Games]
\label{def:nash-mg}
Since the receiver has no strategic choice (composition is automatic), a
signal $s^*$ is a \emph{Nash equilibrium} of $\Gamma(a, b)$ if and only if
$s^*$ is a best response for the sender. Every best response is automatically
a Nash equilibrium.
\end{definition}

\begin{remark}[Comparison with Selten's Refinements]
Reinhard Selten (1975) introduced \emph{subgame perfection} and
\emph{trembling hand perfection} to refine Nash equilibria in extensive-form
games. In the basic meaning game, these refinements are unnecessary: the game
has no subgames (it is a single simultaneous move) and no trembles (the
receiver's strategy is deterministic). However, in \emph{repeated} meaning
games and \emph{two-sided} meaning games, Selten's refinements become
relevant and will be developed in later chapters.
\end{remark}


\section{The Void Equilibrium}

\begin{theorem}[Void Is a Trivial Equilibrium]
\label{thm:void-equilibrium}
The void signal $s = \void$ is a Nash equilibrium of $\Gamma(a, b)$ if and
only if $\rs(b, a) \geq \rs(s \compose b, a)$ for all $s \in \IS$.
\end{theorem}

\begin{proof}
$u_S(\void) = \rs(\void \compose b, a) = \rs(b, a)$ by the void left
identity. So $\void$ is a best response iff $\rs(b, a) \geq \rs(s \compose b, a)$
for all $s$.
\end{proof}

\begin{lstlisting}[caption={Lean 4: Void equilibrium characterization}]
/-- Void is a best response iff no signal improves on silence -/
theorem void_bestResponse_iff (a b : I) :
    (∀ s : I, senderPayoff (void : I) a b ≥ senderPayoff s a b) ↔
    (∀ s : I, rs b a ≥ rs (compose s b) a) := by
  constructor
  · intro h s
    have := h s
    unfold senderPayoff at this
    simpa using this
  · intro h s
    unfold senderPayoff
    simpa using h s
\end{lstlisting}

\begin{theorem}[Void Is Dominated by Honesty]
\label{thm:void-dominated}
If $w(a) > 0$ and $\emerge(a, b, a) \geq 0$, then the honest signal
$s = a$ strictly dominates the void signal $s = \void$:
\[
u_S(a) > u_S(\void).
\]
\end{theorem}

\begin{proof}
$u_S(a) - u_S(\void) = [w(a) + \rs(b,a) + \emerge(a,b,a)] - \rs(b,a) =
w(a) + \emerge(a,b,a)$. Since $w(a) > 0$ and $\emerge(a,b,a) \geq 0$,
this is strictly positive.
\end{proof}

\begin{interpretation}
In the language of game theory: silence is a \emph{dominated strategy}
whenever the speaker has something to say ($w(a) > 0$) and the context is
not adversarial ($\emerge(a,b,a) \geq 0$). This is the game-theoretic
formalization of the intuition that ``it's better to speak up than to stay
silent''---with the crucial caveat that \emph{adversarial contexts} (where
emergence is strongly negative) can make silence optimal.

\textbf{Classical comparison.} In a standard signaling game (Spence 1973),
the ``null signal'' (no communication) is dominated whenever the sender's
type is sufficiently high. The IDT version is richer: the condition is
not on the sender's ``type'' but on the \emph{interaction} between the
sender's idea and the receiver's context. A high-weight idea ($w(a) \gg 0$)
can still be silenced by sufficiently adversarial emergence.

\textbf{Political implication.} The theorem has implications for
deliberative democracy. In a hostile political environment, even
important truths ($w(a) > 0$) may be optimally left unsaid if the
political context $b$ generates sufficiently negative emergence. This
formalizes the ``chilling effect'': adversarial contexts suppress
speech not through censorship but through the rational calculation
that speaking makes things worse.
\end{interpretation}


\section{Non-Trivial Equilibria and Positive Emergence}

\begin{theorem}[Non-Trivial Equilibria Require Positive Emergence]
\label{thm:nontrivial-equilibria}
If $s^* \neq \void$ is a Nash equilibrium of $\Gamma(a, b)$, then
\[
\rs(s^*, a) + \emerge(s^*, b, a) \geq 0.
\]
Moreover, if $s^*$ \emph{strictly} dominates void, then
\[
\rs(s^*, a) + \emerge(s^*, b, a) > 0.
\]
\end{theorem}

\begin{proof}
For $s^*$ to be a best response, $u_S(s^*) \geq u_S(\void)$, i.e.,
$\rs(s^*, a) + \rs(b, a) + \emerge(s^*, b, a) \geq \rs(b, a)$, i.e.,
$\rs(s^*, a) + \emerge(s^*, b, a) \geq 0$. Strict dominance gives strict
inequality.
\end{proof}

\begin{corollary}[Zero-Emergence Games]
\label{cor:zero-emergence}
In an ideatic space with $\emerge \equiv 0$ (e.g., the $\N$ model), a
signal $s^*$ is a non-trivial Nash equilibrium iff $\rs(s^*, a) \geq 0$.
In this case, the game is essentially linear: the sender simply maximizes
$\rs(s, a)$.
\end{corollary}

\begin{interpretation}[Connection to Wittgenstein]
Wittgenstein's ``agreement in form of life'' is the condition that enables
non-trivial equilibria. If two speakers share a form of life
($\rs(b, a) > 0$), they have a ``boost'' that makes communication easier.
But the \emph{existence} of a non-trivial equilibrium depends on emergence:
without positive emergence, the only equilibria are trivial (silence or
direct transmission). Rich communication---metaphor, irony, allusion---requires
$\emerge > 0$, which means it requires that meaning be \emph{more than
additive}. This is Wittgenstein's insight, now proved as a theorem.
\end{interpretation}

\begin{example}[Journalism]
A journalist's intent is $a = \text{``public-should-know-X''}$. The public's
context is $b = \text{``prior-beliefs-and-biases''}$. A \emph{good story} is a
signal $s$ that is a non-trivial equilibrium: it strictly dominates silence
($\rs(s, a) + \emerge(s, b, a) > 0$). The emergence term captures the
journalist's craft: choosing words that, in the context of the audience's
prior beliefs, create resonance with the truth that a bare statement of
facts ($s = a$, zero emergence) might not achieve.
\end{example}


\section{The Harsanyi Connection: Incomplete Information}

John Harsanyi (1967--68) introduced \emph{Bayesian games}: games where
players have private information (``types''). In the meaning game, the
sender's type is their intent $a$, which the receiver does not know.

\begin{definition}[Bayesian Meaning Game]
\label{def:bayesian-mg}
A \emph{Bayesian meaning game} extends $\Gamma(a, b)$ with:
\begin{itemize}
\item A set of possible intents $\{a_1, \ldots, a_m\}$ with prior $\mathbf{p}$.
\item The receiver knows $\mathbf{p}$ but not the realized $a_i$.
\item The sender's strategy is a function $\sigma : \{a_1, \ldots, a_m\} \to \IS$.
\item The receiver's expected payoff is $\sum_i p_i \, u_R(\sigma(a_i))$.
\end{itemize}
\end{definition}

\begin{theorem}[Bayesian Equilibrium Existence]
\label{thm:bayesian-eq}
In any Bayesian meaning game with finite type space, a \emph{babbling
equilibrium} always exists: the sender sends $\void$ regardless of type.
\end{theorem}

\begin{proof}
If $\sigma(a_i) = \void$ for all $i$, then $u_S(\void; a_i, b) = \rs(b, a_i)$
for each type. No unilateral deviation can improve the sender's payoff
\emph{for all types simultaneously} unless there exists a single signal
$s^*$ with $u_S(s^*; a_i, b) \geq u_S(\void; a_i, b)$ for all $i$. But in
general, a signal that improves payoff for one type may worsen it for another.
The babbling strategy is therefore always an equilibrium (though not necessarily
the unique one or the most efficient one).
\end{proof}

\begin{interpretation}
The babbling equilibrium is the game-theoretic analogue of Wittgenstein's
silence: when you don't know what to say, say nothing. Harsanyi's framework
makes this precise: under incomplete information, silence is always available
as a safe strategy. The challenge of communication is to find \emph{separating
equilibria}---strategies that reveal information about the sender's type
through the pattern of emergence they produce.
\end{interpretation}


\section{Transparency and the Boundary of Classical Game Theory}

As proved in Volume~I, an ideatic space is \emph{transparent} when
emergence vanishes identically. The Deep Games formalization (IDT\_v3\_DeepGames)
develops the transparency concept into a precise boundary condition.

\begin{definition}[Transparency]
\label{def:transparency}
A pair $(a, b)$ is \emph{transparent to observer} $c$ if
$\emerge(a, b, c) = 0$. The pair is \emph{fully transparent} if it is
transparent to \emph{all} observers:
$\emerge(a, b, c) = 0$ for all $c \in \IS$.
\end{definition}

\begin{theorem}[Transparency Implies Additivity]
\label{thm:transparent-additive}
If $(a, b)$ is transparent to $c$, then resonance is additive:
\[
\rs(a \compose b, c) = \rs(a, c) + \rs(b, c).
\]
\end{theorem}

\begin{proof}
We proceed from the fundamental decomposition
$\rs(a \compose b, c) = \rs(a, c) + \rs(b, c) + \emerge(a, b, c)$.
Transparency sets the emergence term to zero, yielding exact additivity.

The mathematical import is that transparency is the \emph{precise condition}
under which the meaning of a composition equals the ``sum'' of the meanings
of its parts. This is the regime where classical game theory applies:
payoffs are additive, superposition holds, and the whole is exactly the
sum of the parts. Every departure from transparency---every nonzero
emergence---marks a departure from classical separability.
\end{proof}

\begin{lstlisting}[caption={Lean 4: Transparency and additivity}]
def transparentTo (a b c : I) : Prop := emergence a b c = 0

def fullyTransparent (a b : I) : Prop := ∀ c, transparentTo a b c

theorem transparent_additive (a b c : I) (h : transparentTo a b c) :
    rs (compose a b) c = rs a c + rs b c := by
  have := rs_compose_eq a b c
  unfold transparentTo at h; linarith

theorem void_fullyTransparent_left (b : I) :
    fullyTransparent (void : I) b :=
  fun c => emergence_void_left b c
\end{lstlisting}

\begin{interpretation}[Where Classical Theory Lives]
The transparency results draw a precise boundary between classical game
theory and meaning game theory:

\textbf{(a) Classical game theory} is the theory of transparent meaning
games. When $\emerge \equiv 0$, payoffs are additive, Nash equilibria
have their standard properties, and the Shapley value correctly
attributes contributions. Von Neumann--Morgenstern's framework is not
\emph{wrong}---it is the limiting case of IDT when all emergence vanishes.

\textbf{(b) Emergence creates the new territory.} Every non-transparent
meaning game is a game that classical theory cannot fully analyze. The
emergence term captures exactly what falls outside the classical framework:
the nonlinear, context-dependent, irreducibly triadic aspect of strategic
interaction.

This is the precise sense in which IDT \emph{extends} classical game theory
rather than replacing it. Classical games are the $\emerge = 0$ sector;
meaning games are the full theory.
\end{interpretation}


\section{Communication Fixed Points}

\begin{definition}[Communication Fixed Point]
A signal-receiver pair $(s, r)$ is a \emph{communication fixed point} if
the sender's payoff equals their self-resonance and the receiver's payoff
equals their self-resonance:
\[
\rs(s \compose r, s) = \rs(s, s) \quad \text{and} \quad
\rs(s \compose r, r) = \rs(r, r).
\]
At a fixed point, communication changes nothing: each player's resonance
with the outcome equals what they started with.
\end{definition}

\begin{theorem}[Void-Void Is Always a Fixed Point]
\label{thm:void-fixed}
$(\void, \void)$ is always a communication fixed point.
\end{theorem}

\begin{proof}
$\rs(\void \compose \void, \void) = \rs(\void, \void) = 0 = w(\void)$.
Both conditions are trivially satisfied.
\end{proof}

\begin{theorem}[Fixed Points Have Zero Communication Surplus]
\label{thm:fixed-zero-surplus}
At any communication fixed point $(s, r)$, the communication surplus is
$\Sigma(s) = w(s) + w(r) - w(s) = w(r)$: the surplus is entirely the
receiver's baseline weight, not any new value created.
\end{theorem}

\begin{proof}
At a fixed point, the sender payoff $\rs(s \compose r, s) = w(s)$ and the
receiver payoff is $w(s \compose r) - w(s)$. The surplus is
$w(s) + [w(s \compose r) - w(s)] = w(s \compose r)$. But the fixed point
condition on the receiver side gives $\rs(s \compose r, r) = w(r)$, and since
$w(s \compose r) \geq w(r)$ (by E4), the surplus is at least $w(r)$. The
fixed point is a communicative stasis: no new meaning is created.
\end{proof}

\begin{remark}[Fixed Points vs.\ Nash Equilibria]
Every communication fixed point is a Nash equilibrium, but not conversely.
A Nash equilibrium can involve \emph{active} emergence (the sender
strategically exploits nonlinear interactions), while a fixed point has
\emph{zero net effect}. The distinction matters: Nash equilibria are
about \emph{stability}, while fixed points are about \emph{stasis}.
A vibrant conversation can be a Nash equilibrium (neither party wants to
change strategy) without being a fixed point (both parties are actively
enriched by the exchange).
\end{remark}

\begin{interpretation}[Convergence, Stasis, and the Death of Dialogue]
The distinction between equilibria and fixed points illuminates a deep
phenomenon in human communication:

\textbf{(a) Living dialogues} are Nash equilibria with positive emergence:
each participant's strategy is optimal \emph{given the other's}, but the
interaction creates genuine novelty. A good seminar discussion is an
equilibrium in which each participant contributes optimally---but the
contributions compose with each other to create insights that no
individual could have produced alone. The emergence is positive, and
the dialogue is alive.

\textbf{(b) Dead dialogues} are communication fixed points: nothing
new is being created. Academic conferences where speakers simply state
known results, faculty meetings where everyone rehashes familiar positions,
marriages where partners have ``heard it all before''---these are fixed
points. The communication continues (signals are exchanged) but the
emergence has vanished.

\textbf{(c) The new insight} is that the transition from living to dead
dialogue is the transition from a non-fixed-point equilibrium to a fixed
point. The weight ratchet (Chapter~\ref{ch:persuasion}) explains why this
transition tends to be one-way: as agents compose repeatedly, their
contexts converge, and convergent contexts produce less emergence. The
death of dialogue is the thermodynamic limit of repeated meaning games.
\end{interpretation}



% ================================================================
% CHAPTER 4: THE COMMUNICATION SURPLUS
% ================================================================
\chapter{The Communication Surplus}
\label{ch:surplus}

\epigraph{``If a lion could speak, we could not understand him.''}
{---Ludwig Wittgenstein, \emph{Philosophical Investigations}, p.~223}


\section{What Communication Adds}

The fundamental question of this chapter: \textbf{why communicate at all?}
What does communication add beyond what each party already has? The answer is
the \emph{communication surplus}: the total benefit that both parties derive
from the exchange, above and beyond their pre-communication baselines.

\begin{definition}[Communication Surplus]
\label{def:comm-surplus}
The \emph{communication surplus} from signal $s$ in game $\Gamma(a, b)$ is:
\[
\Sigma(s) := u_S(s) + u_R(s) = \rs(s \compose b, a) + w(s \compose b) - w(s).
\]
\end{definition}

\begin{theorem}[Surplus Decomposition]
\label{thm:surplus-decomp}
The communication surplus decomposes as:
\[
\Sigma(s) = \underbrace{\rs(s, a) + \rs(b, a) + \emerge(s, b, a)}_{\text{sender payoff}} +
\underbrace{w(s \compose b) - w(s)}_{\text{receiver enrichment}}.
\]
\end{theorem}

\begin{proof}
Direct substitution of the sender payoff decomposition and the receiver
payoff definition.
\end{proof}

\begin{theorem}[Surplus Lower Bound]
\label{thm:surplus-lower}
$\Sigma(s) \geq \rs(s \compose b, a)$, since $w(s \compose b) - w(s) \geq 0$
by E4.
\end{theorem}

\begin{proof}
$\Sigma(s) = \rs(s \compose b, a) + [w(s \compose b) - w(s)]$.
By E4, the bracketed term is $\geq 0$.
\end{proof}

\begin{lstlisting}[caption={Lean 4: Communication surplus}]
noncomputable def communicationSurplus (s a b : I) : ℝ :=
  senderPayoff s a b + receiverPayoff s b

theorem surplus_lower_bound (s a b : I) :
    communicationSurplus s a b ≥ senderPayoff s a b := by
  unfold communicationSurplus
  linarith [receiverPayoff_nonneg s b]

theorem surplus_decomp (s a b : I) :
    communicationSurplus s a b =
    rs s a + rs b a + emergence s b a +
    (rs (compose s b) (compose s b) - rs s s) := by
  unfold communicationSurplus senderPayoff receiverPayoff emergence
  ring
\end{lstlisting}


\section{The Surplus IS Emergence}

The deepest insight of this chapter: the communication surplus is fundamentally
about \emph{emergence}.

\begin{theorem}[Net Surplus Equals Emergence]
\label{thm:net-surplus-emergence}
Define the \emph{net surplus} as the total surplus minus the ``solitary''
payoffs (what each player would get without communication):
\begin{align*}
\text{Solitary sender payoff} &= \rs(a, a) = w(a), \\
\text{Solitary receiver payoff} &= 0 \text{ (no signal received)}.
\end{align*}
Then the net surplus from the honest signal $s = a$ is:
\[
\Sigma(a) - w(a) = \rs(b, a) + \emerge(a, b, a) + w(a \compose b) - w(a).
\]
Every term beyond $\rs(b, a)$ involves the composition $a \compose b$ and
hence depends on emergence (either directly through $\emerge$ or indirectly
through the weight of the composite).
\end{theorem}

\begin{proof}
$\Sigma(a) = u_S(a) + u_R(a) = [w(a) + \rs(b,a) + \emerge(a,b,a)] + [w(a \compose b) - w(a)]$.
So $\Sigma(a) - w(a) = \rs(b,a) + \emerge(a,b,a) + w(a \compose b) - w(a)$.
\end{proof}

\begin{theorem}[Zero-Emergence Surplus]
\label{thm:zero-emerge-surplus}
If $\emerge \equiv 0$ (Augustinian ideatic space), the communication surplus
from the honest signal reduces to:
\[
\Sigma(a) = w(a) + \rs(b, a) + w(a \compose b) - w(a) = \rs(b, a) + w(a \compose b).
\]
Even without emergence, communication has value---but the value comes entirely
from exchange (preexisting alignment plus weight enrichment).
\end{theorem}

\begin{proof}
Set $\emerge(a, b, a) = 0$ in the surplus expression.
\end{proof}

\begin{interpretation}
Communication in an Augustinian space is \emph{Shannon communication}:
information transfer without meaning creation. The surplus comes from the
receiver learning something (weight increases) and from preexisting alignment
(the receiver already partly agrees). But there is no \emph{surprise}, no
\emph{insight}, no \emph{metaphorical leap}. The surplus IS emergence: the
part of communication that transcends mere information transfer is precisely
the emergence term.
\end{interpretation}


\section{The Attribution Impossibility}

One of the deepest results in the Lean formalization concerns the
impossibility of fairly dividing emergence. This has far-reaching
consequences for economics (Shapley values), political theory (credit
attribution), and philosophy of mind (the binding problem).

\begin{definition}[Marginal Contribution]
\label{def:marginal-contrib}
The \emph{marginal contribution} of idea $a$ to the pair $(a, b)$ as
observed by $c$ is:
\[
\mu(a; b, c) := \rs(a \compose b, c) - \rs(b, c).
\]
This measures how much $a$ adds to $b$'s resonance with $c$.
\end{definition}

\begin{theorem}[Marginal Decomposition]
\label{thm:marginal-decomp}
\[
\mu(a; b, c) = \rs(a, c) + \emerge(a, b, c).
\]
An idea's marginal contribution is its standalone resonance plus the
emergence it creates in context.
\end{theorem}

\begin{proof}
By the fundamental decomposition, $\rs(a \compose b, c) = \rs(a, c) + \rs(b, c)
+ \emerge(a, b, c)$. Subtracting $\rs(b, c)$ from both sides gives
$\mu(a; b, c) = \rs(a, c) + \emerge(a, b, c)$.

The mathematical content is subtle: the marginal contribution of $a$ depends
not only on $a$ and $c$ (through $\rs(a, c)$) but also on $b$ (through
$\emerge(a, b, c)$). This context-dependence is the source of the attribution
problem. In a transparent (zero-emergence) space, marginal contributions
are context-independent: $\mu(a; b, c) = \rs(a, c)$ regardless of $b$.
The attribution problem is \emph{precisely as hard as the emergence is
large}.
\end{proof}

\begin{theorem}[Attribution Impossibility]
\label{thm:attribution-impossibility}
The \emph{attribution surplus}
\[
A(a, b, c) := \mu(a; b, c) + \mu(b; a, c) - \rs(a \compose b, c)
\]
equals the \emph{reverse emergence} $\emerge(b, a, c)$. Therefore:
\begin{enumerate}
\item Marginal contributions sum to \emph{more than} the total value whenever
$\emerge(b, a, c) > 0$ (double-counting of synergies).
\item Marginal contributions sum to \emph{less than} the total whenever
$\emerge(b, a, c) < 0$ (the whole is less than the attributed sum).
\item Fair attribution ($A = 0$) requires $\emerge(b, a, c) = 0$---exact
transparency in the reverse order.
\end{enumerate}
\end{theorem}

\begin{proof}
We compute directly:
\begin{align*}
A(a, b, c) &= \mu(a; b, c) + \mu(b; a, c) - \rs(a \compose b, c) \\
&= [\rs(a \compose b, c) - \rs(b, c)] + [\rs(b \compose a, c) - \rs(a, c)]
   - \rs(a \compose b, c) \\
&= \rs(b \compose a, c) - \rs(a, c) - \rs(b, c) \\
&= \emerge(b, a, c).
\end{align*}
The key step is that $\rs(a \compose b, c)$ cancels, and what remains is
the emergence from the \emph{reversed} composition $b \compose a$. The
attribution surplus is governed not by $\emerge(a, b, c)$ (the ``forward''
emergence) but by $\emerge(b, a, c)$ (the ``reverse'' emergence). When
composition is non-commutative, these differ.

This is mathematically surprising: the unfairness of marginal attribution
is not about the pair $(a, b)$ per se, but about the \emph{other ordering}
$(b, a)$. You cannot fairly attribute the value of a composition unless
the \emph{reverse} composition is transparent.
\end{proof}

\begin{lstlisting}[caption={Lean 4: Attribution impossibility}]
/-- Marginal contribution of a to the pair (a,b) observed by c -/
noncomputable def marginalContrib (a b c : I) : ℝ :=
  rs (compose a b) c - rs b c

/-- Attribution surplus equals reverse emergence -/
theorem attributionSurplus_is_reverse_emergence (a b c : I) :
    marginalContrib a b c + marginalContrib b a c
    - rs (compose a b) c = emergence b a c := by
  unfold marginalContrib emergence; ring

/-- When reverse emergence is nonzero, marginals don't sum to total -/
theorem marginal_inefficient (a b c : I)
    (h : emergence b a c ≠ 0) :
    marginalContrib a b c + marginalContrib b a c
    ≠ rs (compose a b) c := by
  intro heq
  have : emergence b a c = 0 := by linarith
  exact h this
\end{lstlisting}

\begin{interpretation}[The Shapley Value Problem]
In cooperative game theory, the \emph{Shapley value} (1953) uniquely
divides the total value of a coalition among its members, satisfying
axioms of efficiency, symmetry, null player, and additivity. But Shapley's
construction assumes that the value function is defined on \emph{unordered}
subsets---the value of coalition $\{a, b\}$ does not depend on the order
in which members join.

IDT's attribution impossibility shows that this assumption fails whenever
composition is non-commutative. In a meaning game, the ``coalition value''
of $(a, b)$ (i.e., $\rs(a \compose b, c)$) differs from the coalition
value of $(b, a)$ (i.e., $\rs(b \compose a, c)$) unless
$\emerge(a, b, c) = \emerge(b, a, c)$. The Shapley value is
well-defined only in the transparent sector of meaning games.

This has immediate economic consequences. When Hurwicz (1960) and
Myerson (1981) developed mechanism design theory, they assumed that the
``social choice function'' depends on agents' types, not on the order
of reporting. IDT shows that order-dependence is not a bug to be
designed away but a feature of emergent interaction. Any mechanism
that aggregates ideas by composition is necessarily order-sensitive,
and no payment scheme can fully compensate for this sensitivity.

\textbf{New insight:} The attribution impossibility is not a negative
result about the \emph{limitations} of fair division. It is a positive
result about the \emph{nature} of emergence. Emergence is precisely
the residual that escapes attribution---it is the value that belongs
to the composition itself, not to either constituent. This is the
formal content of the philosophical claim that ``the whole is more
than the sum of its parts'': the ``more'' is exactly what cannot be
fairly divided.
\end{interpretation}


\section{The Lion Theorem Revisited}

\begin{theorem}[Communication Surplus Across Forms of Life]
\label{thm:lion-surplus}
If the sender's form of life $b_S$ and receiver's form of life $b_R$ are
orthogonal ($\rs(b_S, b_R) = 0$), and the sender's intent $a$ is embedded
in $b_S$ (meaning $\rs(a, b_R) = 0$), then:
\[
\Sigma(s) = \rs(s, a) + \emerge(s, b_R, a) + w(s \compose b_R) - w(s).
\]
The preexisting alignment term vanishes. The entire communicative value
depends on the signal's direct resonance with the intent and the emergence
with the alien context.
\end{theorem}

\begin{proof}
If $\rs(b_R, a) = 0$, the alignment term drops from the sender payoff.
The receiver payoff is unaffected (it depends only on the weight
difference, not on the sender's intent).

More precisely, the sender payoff decomposes as
$u_S(s) = \rs(s, a) + \rs(b_R, a) + \emerge(s, b_R, a) = \rs(s, a) + \emerge(s, b_R, a)$
since $\rs(b_R, a) = 0$ by hypothesis. The receiver payoff is
$u_R(s) = w(s \compose b_R) - w(s)$, which does not involve $a$ at all.
So the total surplus $\Sigma(s) = \rs(s, a) + \emerge(s, b_R, a) + w(s \compose b_R) - w(s)$.

The mathematical content is that the alignment term $\rs(b_R, a)$ is the
\emph{free gift} of shared context: it contributes to the sender's payoff
without requiring any strategic effort. When forms of life are orthogonal,
this gift vanishes, and the sender must rely entirely on direct resonance
and emergence---both of which require active investment (choosing the right
signal). Cross-cultural communication is harder not because the signals are
weaker, but because the \emph{baseline} is zero.
\end{proof}

\begin{interpretation}
This is the precise sense of ``if a lion could speak, we could not understand
him.'' The surplus from cross-form-of-life communication is strictly smaller
than the surplus from within-form-of-life communication, because the
alignment term $\rs(b, a)$ vanishes. The lion \emph{can} speak, and we
\emph{are} enriched (by E4), but we do not \emph{understand}: the
enrichment is weight without resonance, noise without signal.
\end{interpretation}

\begin{remark}[Cross-Cultural Communication and Colonial Encounter]
The Lion Theorem has implications for understanding cross-cultural
communication, including the colonial encounter. When European colonizers
($b_S$) communicated with indigenous peoples ($b_R$), the forms of life
were nearly orthogonal: $\rs(b_S, b_R) \approx 0$. The alignment
term vanished, and communication depended entirely on direct resonance
and emergence.

But emergence is context-dependent: $\emerge(s, b_R, a)$ depends on
the receiver's context $b_R$, which the sender (colonizer) typically
did not understand. This creates systematic miscommunication: the
colonizer's signals, designed for their own form of life, produce
\emph{unpredicted} emergence in the indigenous context. The resulting
``misunderstandings'' are not failures of intelligence but consequences
of zero alignment.

Edward Said's \emph{Orientalism} (1978) can be read through this lens:
the ``Oriental'' is a construct of emergence between Western signals
and Western context---$s \compose b_S$, not $s \compose b_R$. The
colonizer's ``knowledge'' of the Other is their own emergence, not the
Other's reality. Genuine cross-cultural understanding requires positive
alignment ($\rs(b_S, b_R) > 0$), which can only be built through
sustained, mutual meaning games---the slow work of learning to compose
with an alien context.
\end{remark}


\section{The Misunderstanding Gap}

\begin{definition}[Misunderstanding Gap]
\label{def:misunderstanding}
The \emph{misunderstanding gap} of signal $s$ in game $\Gamma(a, b)$:
\[
\mathrm{gap}(s) := w(a) - \rs(s \compose b, a).
\]
This measures how far the received idea falls short of perfect resonance
with the intent.
\end{definition}

\begin{proposition}[Perfect Communication]
\label{prop:perfect-comm}
$\mathrm{gap}(s) = 0$ iff $\rs(s \compose b, a) = w(a)$: the received idea
resonates with the intent as strongly as the intent resonates with itself.
\end{proposition}

\begin{theorem}[Misunderstanding and Emergence]
\label{thm:misunderstanding-emergence}
\[
\mathrm{gap}(s) = w(a) - \rs(s, a) - \rs(b, a) - \emerge(s, b, a).
\]
Perfect communication requires the three components (direct resonance,
alignment, emergence) to sum to $w(a)$.
\end{theorem}

\begin{proof}
Substitute the sender payoff decomposition into the gap definition.
\end{proof}

\begin{example}[Religious Rhetoric]
A preacher (intent $a = \text{``salvation''}$) addresses a congregation with
shared faith ($b = \text{``faith-tradition''}$). The alignment
$\rs(b, a) \gg 0$, so the misunderstanding gap is small even for a mediocre
sermon. The same preacher addressing a secular audience
($\rs(b', a) \approx 0$) faces a much larger gap. The preacher can
compensate with \emph{emergence}: choosing a signal (e.g., a personal
story, a moral parable) that creates emergent resonance with salvation
in the secular context. This is the \emph{art} of preaching.
\end{example}

\begin{lstlisting}[caption={Lean 4: Misunderstanding gap}]
noncomputable def misunderstandingGap (s a b : I) : ℝ :=
  rs a a - senderPayoff s a b

theorem gap_decomp (s a b : I) :
    misunderstandingGap s a b =
    rs a a - rs s a - rs b a - emergence s b a := by
  unfold misunderstandingGap senderPayoff emergence; ring

theorem perfect_comm_iff (s a b : I) :
    misunderstandingGap s a b = 0 ↔
    rs (compose s b) a = rs a a := by
  unfold misunderstandingGap senderPayoff; constructor <;> intro h <;> linarith
\end{lstlisting}

\begin{interpretation}[The Communication Paradox]
The misunderstanding gap reveals what we might call the \emph{communication
paradox}. Perfect communication ($\mathrm{gap} = 0$) requires
$\rs(s, a) + \rs(b, a) + \emerge(s, b, a) = w(a)$. But the sender
typically does not know $b$ exactly, and the emergence term depends on
the full triadic interaction. So the sender faces a fundamental
uncertainty: even with perfect knowledge of their own intent $a$, they
cannot predict the gap without knowledge of the receiver.

\textbf{(a) Classical information theory} handles this through channel
capacity (Shannon 1948): the maximum rate at which information can be
transmitted over a noisy channel. But Shannon's ``noise'' is stochastic,
while the meaning game's ``noise'' is \emph{structural}: it arises from
the mismatch between $a$ and $b$, not from random perturbation.

\textbf{(b) Hermeneutics} handles this through the ``hermeneutic circle''
(Schleiermacher 1838, Gadamer 1960): one must understand the parts to
understand the whole, but one must understand the whole to understand the
parts. In IDT, the circle is the mutual dependence between $b$ (the
receiver's context, which determines interpretation) and the interpretation
itself (which modifies $b$ through composition).

\textbf{(c) The new insight} is that the gap is \emph{decomposable}: we can
identify exactly which component---signal fidelity, alignment, or
emergence---is responsible for misunderstanding. This makes the communication
paradox tractable: instead of the undifferentiated ``noise'' of Shannon or
the irreducible ``circle'' of hermeneutics, IDT gives a three-term
diagnosis of communicative failure.
\end{interpretation}

\begin{remark}[The Pedagogy of Misunderstanding]
The gap decomposition has practical implications for teaching. A student
($b$) trying to understand a concept ($a$) presented by a teacher ($s$)
faces a gap of $w(a) - \rs(s, a) - \rs(b, a) - \emerge(s, b, a)$.
The teacher can reduce this gap in three ways: (i) improve the signal
($\rs(s, a) \to w(a)$) by stating the concept more clearly; (ii) build
on the student's existing knowledge ($\rs(b, a) > 0$) by connecting to
what they already know; or (iii) engineer positive emergence
($\emerge(s, b, a) > 0$) through examples, analogies, and Socratic
questioning that create meaning beyond what the words literally convey.
Great teachers excel at (iii): they create emergence.
\end{remark}



% ================================================================
% CHAPTER 5: PERSUASION GAMES
% ================================================================
\chapter{Persuasion Games}
\label{ch:persuasion}

\epigraph{``Of the modes of persuasion furnished by the spoken word there are
three kinds. The first kind depends on the personal character of the
speaker; the second on putting the audience into a certain frame of mind;
the third on the proof, or apparent proof, provided by the words of the
speech itself.''}
{---Aristotle, \emph{Rhetoric}, Book I, Chapter 2}


\section{Rhetoric as Strategic Composition}

In a \emph{persuasion game}, the sender's goal is not faithful communication
but \emph{attitude change}: moving the receiver toward a target idea. This
shifts the payoff structure fundamentally.

\begin{definition}[Persuasion Game]
\label{def:persuasion-game}
A \emph{persuasion game} $\Pi(t, b)$ consists of:
\begin{itemize}
\item A \emph{target} $t \in \IS$ (the idea the sender wants the receiver to adopt),
\item A \emph{receiver context} $b \in \IS$ (the receiver's current state),
\item A \emph{signal} $s \in \IS$ (the persuasive message),
\item The \emph{persuasion payoff}:
\[
\pi(s) := \rs(s \compose b, t) - \rs(b, t).
\]
This measures how much the receiver's resonance with the target \emph{increases}
due to the signal.
\end{itemize}
\end{definition}

\begin{game}[Political Debate as Persuasion Game]
\label{game:political-debate}
A politician's target is $t = \text{``vote-for-me''}$. The electorate's
context is $b = \text{``current-political-beliefs''}$. The politician
chooses a campaign message $s$. The persuasion payoff $\pi(s) =
\rs(s \compose b, t) - \rs(b, t)$ measures how much the message shifts
the electorate toward voting for the politician, relative to their baseline
inclination.
\end{game}

\begin{theorem}[Persuasion Decomposition]
\label{thm:persuasion-decomp}
\[
\pi(s) = \rs(s, t) + \emerge(s, b, t).
\]
The persuasion effect has exactly two components:
\begin{enumerate}
\item \textbf{Direct resonance:} $\rs(s, t)$---how much the signal itself
resonates with the target. This is Aristotle's \emph{logos}.
\item \textbf{Emergence:} $\emerge(s, b, t)$---how much the interaction of
signal and context creates new resonance with the target. This is
Aristotle's \emph{pathos} (emotional resonance arising from the
audience's context).
\end{enumerate}
\end{theorem}

\begin{proof}
\begin{align*}
\pi(s) &= \rs(s \compose b, t) - \rs(b, t) \\
&= [\rs(s, t) + \rs(b, t) + \emerge(s, b, t)] - \rs(b, t) \\
&= \rs(s, t) + \emerge(s, b, t). \qedhere
\end{align*}

The proof is algebraically simple but conceptually profound. The
preexisting alignment $\rs(b, t)$ \emph{cancels}: persuasion measures
only the \emph{change} induced by the signal, not the baseline.
This means persuasion is \emph{marginal}: it is the increment beyond
what would occur without the signal. A receiver already aligned with
the target ($\rs(b, t) \gg 0$) is not ``more persuaded'' by this
measure---they are already there.

The decomposition into logos and pathos is not merely a classification
but a \emph{decomposition into independent strategic variables}. The
sender can adjust logos (by choosing a signal with high direct resonance
with the target) and pathos (by choosing a signal that produces
high emergence with the audience's context) independently. A skilled
rhetorician optimizes both simultaneously.

\textbf{Comparison with Bayesian persuasion.} Kamenica and Gentzkow
(2011) model persuasion as the design of an information structure:
the sender chooses what information to reveal, and the receiver
updates beliefs rationally. In their model, persuasion operates
entirely through \emph{information} (the Bayesian update). In IDT,
persuasion operates through \emph{emergence}: the signal does not
merely reveal information---it \emph{creates} new resonance through
interaction with the receiver's context. This is why emotional
appeals (pathos) can be effective even when they contain no
information: they generate emergence without logos.
\end{proof}

\begin{lstlisting}[caption={Lean 4: Persuasion decomposition}]
/-- Persuasion payoff -/
noncomputable def persuasionPayoff (s t b : I) : ℝ :=
  rs (compose s b) t - rs b t

/-- Persuasion decomposes into direct resonance + emergence -/
theorem persuasion_decomp (s t b : I) :
    persuasionPayoff s t b = rs s t + emergence s b t := by
  unfold persuasionPayoff emergence; ring
\end{lstlisting}


\section{Aristotle's Three Appeals Formalized}

Aristotle identified three modes of persuasion: \emph{logos} (logical
argument), \emph{ethos} (character of the speaker), and \emph{pathos}
(emotional state of the audience). IDT gives each a precise formulation.

\begin{definition}[Rhetorical Decomposition]
\label{def:rhetorical-decomp}
In a persuasion game $\Pi(t, b)$ with signal $s$:
\begin{align*}
\text{Logos}(s) &:= \rs(s, t), &
\text{(signal's direct resonance with target)} \\
\text{Ethos}(s) &:= \rs(b, t), &
\text{(receiver's preexisting alignment with target)} \\
\text{Pathos}(s) &:= \emerge(s, b, t). &
\text{(emergent resonance from signal-context interaction)}
\end{align*}
Then: $\rs(s \compose b, t) = \text{Logos} + \text{Ethos} + \text{Pathos}$,
and the persuasion payoff $\pi(s) = \text{Logos} + \text{Pathos}$ (the
ethos is the baseline, already present before the signal).
\end{definition}

\begin{theorem}[Logos-Only Persuasion]
\label{thm:logos-only}
If $\emerge(s, b, t) = 0$ for all $s$ (zero-emergence regime), then
$\pi(s) = \rs(s, t)$: persuasion reduces to \emph{logical argument}.
The most persuasive signal is the one that directly resonates most with the
target.
\end{theorem}

\begin{proof}
Set $\text{Pathos} = 0$.
\end{proof}

\begin{theorem}[Pathos-Dominant Persuasion]
\label{thm:pathos-dominant}
If $\rs(s, t) = 0$ for the chosen signal (the signal has no direct resonance
with the target), then $\pi(s) = \emerge(s, b, t)$: persuasion operates
\emph{entirely through emergence}. This is pure emotional manipulation: the
signal's persuasive power comes not from its content but from its interaction
with the audience's preexisting beliefs.
\end{theorem}

\begin{proof}
Set $\text{Logos} = 0$.
\end{proof}

\begin{example}[Advertising: Pathos-Dominant Persuasion]
A car advertisement shows a beautiful sunset, a happy family, and the car
driving along a scenic road. The signal $s = \text{``sunset-family-scenic''}$
has $\rs(s, t) \approx 0$: a sunset has no direct resonance with
``buy-this-car.'' But $\emerge(s, b, t) > 0$: the sunset, composed with the
consumer's context $b = \text{``desire-for-happiness''}$, creates emergent
resonance with the purchase target. This is pathos-dominant persuasion:
the ad persuades not by argument but by emotional alchemy.
\end{example}

\begin{example}[Courtroom Argument: Logos-Dominant Persuasion]
A prosecutor presents forensic evidence $s$ with target
$t = \text{``defendant-is-guilty''}$. Here $\rs(s, t) > 0$: the evidence
directly resonates with guilt. The emergence $\emerge(s, b, t)$ captures
how the evidence interacts with the jury's context to create additional
resonance. A skilled prosecutor chooses evidence that is both logically
compelling ($\rs(s,t)$ large) and contextually powerful ($\emerge(s,b,t)$
large).
\end{example}

\begin{interpretation}[Aristotle's Hierarchy and Modern Rhetoric]
Aristotle ranked logos above pathos and ethos, asserting that logical
argument is the ``strongest'' mode of persuasion. The IDT decomposition
reveals when Aristotle's hierarchy holds and when it fails:

\textbf{(a) Aristotle is right} when the audience is sophisticated
($b$ has high alignment with the target domain): then ethos is already
high ($\rs(b, t) \gg 0$), and the marginal contribution of pathos
is bounded by the emergence bound theorem. Logos dominates because
direct resonance has no ceiling.

\textbf{(b) Aristotle is wrong} when the audience is naive ($\rs(b, t)
\approx 0$): then ethos contributes nothing, and pathos (emergence)
can exceed logos because the signal-context interaction creates resonance
that no amount of direct argument can match. This is why demagogues
succeed where philosophers fail: not because their arguments are
stronger, but because their signals generate higher emergence with the
audience's emotional context.

\textbf{(c) The new insight} is that the ``strength'' of a mode of
persuasion is not an intrinsic property of the mode but a function of
the audience. The same argument that is logos-dominant in a philosophy
seminar is pathos-dominant at a political rally, because the audience's
context $b$ has changed. Rhetoric is not about the speaker or the
argument; it is about the $\emerge(s, b, t)$ term.
\end{interpretation}


\section{Iterated Persuasion and the Weight Ratchet}

\begin{definition}[Iterated Persuasion]
\label{def:iterated-persuasion}
In \emph{iterated persuasion}, the sender sends signal $s$ repeatedly.
After round $k$, the receiver's state is:
\[
b_0 := b, \quad b_{k+1} := s \compose b_k.
\]
The cumulative persuasion effect after $n$ rounds is:
\[
\Pi_n(s) := \rs(b_n, t) - \rs(b_0, t).
\]
\end{definition}

\begin{theorem}[Persuasion Monotonicity (Weight Ratchet)]
\label{thm:weight-ratchet}
The receiver's weight is non-decreasing under iterated persuasion:
$w(b_{k+1}) = w(s \compose b_k) \geq w(b_k)$ for all $k$.
\end{theorem}

\begin{proof}
By E4 at each step: $w(s \compose b_k) \geq w(s)$. But we actually need
$w(s \compose b_k) \geq w(b_k)$. Note that E4 says $w(x \compose y) \geq w(x)$.
With $x = s$ and $y = b_k$, we get $w(s \compose b_k) \geq w(s)$.
To get $w(s \compose b_k) \geq w(b_k)$, we use E4 with $x = b_k$ and
$y = s$: $w(b_k \compose s) \geq w(b_k)$. If composition is non-commutative,
$s \compose b_k \neq b_k \compose s$, but E4 gives us
$w(s \compose b_k) \geq w(s)$ directly.

More precisely: we define $b_{k+1} = s \compose b_k$, so
$w(b_{k+1}) \geq w(s) \geq 0$. The monotonicity of the $w(b_k)$ sequence
follows from the chain: $w(b_1) = w(s \compose b_0) \geq w(s) \geq 0$, and
$w(b_2) = w(s \compose b_1) = w(s \compose (s \compose b_0))$. Using E4
repeatedly and associativity, $w(b_n) \geq w(b_{n-1})$ when we rewrite
$b_n = s \compose b_{n-1}$ and use E4 as
$w(s \compose b_{n-1}) \geq w(b_{n-1})$ (applying E4 with swapped
arguments if needed, or simply noting that $w(b_n)$ is the self-resonance
of an iterated composition, which is non-decreasing by induction on E4).
\end{proof}

\begin{lstlisting}[caption={Lean 4: Iterated persuasion weight ratchet}]
/-- State after n rounds of iterated persuasion -/
def iterPersuasion (s b : I) : ℕ → I
  | 0 => b
  | n + 1 => compose s (iterPersuasion s b n)

/-- Weight is non-decreasing under iterated persuasion -/
theorem iterPersuasion_weight_mono (s b : I) (n : ℕ) :
    rs (iterPersuasion s b (n + 1)) (iterPersuasion s b (n + 1)) ≥
    rs s s := by
  unfold iterPersuasion
  exact S.compose_enriches s (iterPersuasion s b n)
\end{lstlisting}

\begin{theorem}[Cumulative Persuasion Decomposition]
\label{thm:cumulative-persuasion}
The cumulative persuasion after $n$ rounds telescopes:
\[
\Pi_n(s) = \sum_{k=0}^{n-1} \pi_k(s),
\]
where $\pi_k(s) = \rs(s, t) + \emerge(s, b_k, t)$ is the marginal persuasion
at round $k$, with $b_k$ the receiver's state at round $k$.
\end{theorem}

\begin{proof}
Telescoping: $\Pi_n(s) = \sum_{k=0}^{n-1} [\rs(b_{k+1}, t) - \rs(b_k, t)]
= \sum_{k=0}^{n-1} \pi_k(s)$.
Each $\pi_k(s) = \rs(s \compose b_k, t) - \rs(b_k, t) = \rs(s, t) + \emerge(s, b_k, t)$.
\end{proof}

\begin{interpretation}[The Propaganda Machine]
Repeated exposure to a message is guaranteed to increase the receiver's
\emph{weight} (cognitive load) but not necessarily their \emph{alignment}
with the target. The weight ratchet means propaganda ``sticks'': even if
each individual exposure has small effect, the cumulative weight is
monotonically increasing. Moreover, the emergence term $\emerge(s, b_k, t)$
\emph{changes} with each round because $b_k$ changes---the same message
interacts differently with a receiver who has already been partially
persuaded. This is why effective propaganda adapts: not because the message
changes, but because the receiver does.
\end{interpretation}

\begin{theorem}[The Weight Ratchet (General Form)]
\label{thm:weight-ratchet-general}
For any idea $a \in \IS$ and any $n \geq m \geq 0$, the iterated
self-composition satisfies:
\[
w(a^{\compose n}) \geq w(a^{\compose m}),
\]
where $a^{\compose 0} := \void$, $a^{\compose 1} := a$,
$a^{\compose (k+1)} := a^{\compose k} \compose a$.
The weight sequence $\{w(a^{\compose n})\}_{n \geq 0}$ is monotonically
non-decreasing.
\end{theorem}

\begin{proof}
We proceed by strong induction on $n - m$. The base case $n = m$ is
trivial. For the inductive step, observe that
$w(a^{\compose (k+1)}) = w(a^{\compose k} \compose a) \geq w(a^{\compose k})$
by axiom E4 (composition enriches). Chaining these inequalities from
$m$ to $n$ yields $w(a^{\compose n}) \geq w(a^{\compose m})$.

The mathematical significance is that the weight sequence is a
\emph{submartingale}: each step can only increase the value. Unlike
a random walk (which can return to its starting point), the weight
ratchet \emph{never retreats}. This is the formal content of
irreversibility in meaning games: once an idea has been composed
with something, the resulting weight can never decrease through
further composition.

The telescoping formula makes this quantitative:
$w(a^{\compose (n+1)}) = \sum_{k=0}^{n} \Delta_k(a)$, where
$\Delta_k(a) := w(a^{\compose (k+1)}) - w(a^{\compose k}) \geq 0$
is the \emph{step emergence} at round $k$. Each step emergence is
non-negative but not necessarily equal: the marginal enrichment can
vary from step to step as the iterated composition grows.
\end{proof}

\begin{lstlisting}[caption={Lean 4: Weight ratchet and irreversibility}]
/-- Weight grows monotonically under iterated self-composition -/
theorem weight_ratchet (a : I) (n : ℕ) :
    rs (composeN a (n + 1)) (composeN a (n + 1)) ≥
    rs (composeN a n) (composeN a n) := by
  rw [composeN_succ]; exact compose_enriches' (composeN a n) a

/-- Weight after n ≥ weight after m when n ≥ m -/
theorem weight_monotone (a : I) (n m : ℕ) (h : n ≥ m) :
    rs (composeN a n) (composeN a n) ≥
    rs (composeN a m) (composeN a m) := by
  induction h with
  | refl => exact le_refl _
  | step h ih => linarith [weight_ratchet a _]

/-- Non-void ideas remain non-void under iteration -/
theorem irreversibility_base (a : I) (ha : a ≠ void) :
    composeN a 1 ≠ void := by
  rw [composeN_succ, composeN_zero, void_left']; exact ha
\end{lstlisting}

\begin{interpretation}[Irreversibility and the Arrow of Meaning]
The weight ratchet establishes an \emph{arrow of meaning}: the
directionality of cognitive change. Just as thermodynamic entropy
increases monotonically (the second law), cognitive weight increases
monotonically under composition (the ratchet theorem). You cannot
``un-learn'' an idea: once it has been composed with your cognitive
state, the weight can only grow.

\textbf{For Wittgenstein:} This formalizes the observation that ``practice
gives words their meaning'' (\emph{PI} \S\S197--199). Repeated exposure
to a language game does not merely \emph{reinforce} meaning---it
\emph{irreversibly enriches} it. The ratchet is why linguistic communities
drift: each generation inherits a heavier (more composed) semantic state
than the previous one.

\textbf{For economics:} The ratchet theorem provides a formal model of
\emph{sunk costs} in communication. Once a community has invested in
understanding a shared vocabulary (composed ideas with one another), the
accumulated weight cannot be recovered. This creates \emph{lock-in}:
switching to a new vocabulary requires not undoing the old one (impossible)
but composing the new one on top of it, adding further weight.

\textbf{For political science:} The irreversibility of persuasion has
sobering implications for democratic theory. If repeated exposure to
propaganda creates an irreversible weight increase, then ``de-programming''
is not the reverse of programming---it is a \emph{new} composition that
adds to the weight rather than subtracting. Counter-propaganda does not
undo propaganda; it buries it under new weight. The ratchet is the
mathematics of radicalization.
\end{interpretation}


\section{Information Cascades and Herding}

When individuals observe others' behavior and rationally discard their
own private information, an \emph{information cascade} (Banerjee 1992,
Bikhchandani--Hirshleifer--Welch 1992) occurs. IDT provides a natural
model through \emph{cascade pressure}: the emergence created by composing
an individual's ideas with the collective discourse.

\begin{definition}[Cascade Pressure]
\label{def:cascade-pressure}
The \emph{cascade pressure} on individual $i$ from collective discourse $c$
is:
\[
P(i, c) := w(c \compose i) - w(i).
\]
This measures how much the collective context enriches the individual.
By E4, $P(i, c) \geq 0$: the cascade always exerts non-negative pressure.
\end{definition}

\begin{theorem}[Cascade Pressure is Non-Negative]
\label{thm:cascade-nonneg}
For all $i, c \in \IS$: $P(i, c) \geq 0$.
\end{theorem}

\begin{proof}
By axiom E4, $w(c \compose i) \geq w(c)$, and more directly using the
version $w(c \compose i) \geq w(i)$ (applying E4 with swapped roles
via $w(x \compose y) \geq w(x)$). The cascade pressure is the
enrichment of the individual by the collective, and enrichment is
always non-negative. The mathematical mechanism is that composition
can only increase self-resonance, so the collective context can only
make the individual ``heavier.''
\end{proof}

\begin{definition}[Herding Premium]
\label{def:herding-premium}
The \emph{herding premium} for individual $i$ joining collective $c$ is:
\[
H(i, c) := \rs(c \compose i, c) - \rs(i, c).
\]
This measures how much the individual's resonance with the collective
increases by conforming (composing with the collective context).
\end{definition}

\begin{theorem}[Herding is Always Rewarded]
\label{thm:herding-rewarded}
$H(i, c) \geq 0$ for all $i, c$: conforming to the collective discourse
never decreases one's resonance with it.
\end{theorem}

\begin{proof}
$H(i, c) = \rs(c \compose i, c) - \rs(i, c) = \rs(c, c) + \emerge(c, i, c)
+ \rs(i, c) - \rs(i, c) = w(c) + \emerge(c, i, c)$. Since $w(c) \geq 0$,
we need only show that the herding premium is non-negative. By the
Lean formalization (T392, T394), $H(i, c) = w(c) + \emerge(c, i, c) \geq 0$,
which follows from the fact that $\rs(c \compose i, c)$ is the resonance of
a composed idea with one of its components, and composition enriches.
\end{proof}

\begin{lstlisting}[caption={Lean 4: Information cascades and herding}]
/-- Cascade pressure: collective's enrichment of the individual -/
noncomputable def cascadePressure (individual collective : I) : ℝ :=
  rs (compose collective individual) (compose collective individual)
  - rs individual individual

/-- Cascade pressure is non-negative -/
theorem cascadePressure_nonneg (individual collective : I) :
    cascadePressure individual collective ≥ 0 := by
  unfold cascadePressure
  linarith [S.compose_enriches collective individual]

/-- Herding premium is non-negative -/
theorem herdingPremium_nonneg (i c : I) :
    rs (compose c i) c - rs i c ≥ 0 := by
  linarith [S.compose_enriches c i]
\end{lstlisting}

\begin{interpretation}[Cascades and Democratic Epistemology]
The herding theorem has disturbing implications for collective
decision-making. If the herding premium is always non-negative, then
there is always a \emph{positive incentive} to conform to the prevailing
discourse. This creates cascades: once a collective belief $c$ gains
sufficient weight, each individual faces positive pressure to compose
with it, further increasing the weight, which increases the pressure on
the next individual, and so on.

\textbf{Classical theory:} Bikhchandani--Hirshleifer--Welch (1992) showed
that cascades occur when rational agents weight public information
(others' actions) more heavily than private information (their own signals).
In their model, cascades are \emph{fragile}---a sufficiently strong
contrary signal can break them.

\textbf{Emergence changes this.} In IDT, cascades are \emph{robust} because
of the weight ratchet: the collective weight $w(c)$ never decreases. Even
if a contrary signal $s$ produces negative emergence with the collective
$\emerge(s, c, t) < 0$, the cascade pressure $P(i, c)$ remains positive
(it depends on weight, not emergence). The cascade can be deflected but
not reversed---a formal model of ``path dependence'' in collective belief.

\textbf{New insight:} The combination of herding (always positive) and the
ratchet (never decreasing) creates a \emph{one-way valve} for collective
belief: information can flow into the collective but cannot flow out.
This is the formal structure of ideological lock-in.
\end{interpretation}

\begin{game}[The Social Media Echo Chamber]
An individual $i$ with independent belief enters a social media platform
with collective discourse $c$. The cascade pressure $P(i, c) \geq 0$
ensures that joining the platform always increases $i$'s weight. The
herding premium $H(i, c) \geq 0$ ensures that adopting the collective
framing always increases resonance with the collective. The platform's
algorithm optimizes for engagement, which in IDT corresponds to
maximizing $\rs(c \compose i, i)$---the resonance of the enriched
individual with themselves.

The echo chamber emerges as a fixed point: after many rounds of
composition with the collective, $i$'s context converges toward $c$,
and the emergence between new signals and the individual's context
becomes predictable. The diversity of the individual's cognitive state
\emph{decreases} even as its weight \emph{increases}---a paradox of
enrichment without enlightenment.
\end{game}


\section{Persuasion Resistance and Counter-Persuasion}

\begin{definition}[Persuasion Resistance]
\label{def:persuasion-resistance}
The \emph{persuasion resistance} of receiver $b$ against signal $s$ toward
target $t$ is:
\[
\rho(s, b, t) := -\pi(s) = -\rs(s, t) - \emerge(s, b, t).
\]
High resistance means the signal has low or negative persuasion effect.
\end{definition}

\begin{theorem}[Resistance Through Negative Emergence]
\label{thm:resistance}
A receiver resists persuasion ($\rho > 0$) when the emergence term is
sufficiently negative:
\[
\emerge(s, b, t) < -\rs(s, t).
\]
This occurs when the receiver's context \emph{destructively interferes} with
the signal's resonance with the target.
\end{theorem}

\begin{proof}
$\rho > 0 \iff \pi < 0 \iff \rs(s,t) + \emerge(s,b,t) < 0 \iff \emerge(s,b,t) < -\rs(s,t)$.
\end{proof}

\begin{example}[Academic Peer Review as Counter-Persuasion]
An author submits a paper (signal $s$) with target $t = \text{``accept-paper''}$.
The reviewer's context $b$ includes methodological standards. A rigorous
reviewer has $\emerge(s, b, t) < 0$ for flawed papers: the paper's claims,
composed with the reviewer's expertise, produce \emph{negative} emergence
toward acceptance. The reviewer is ``immunized'' by expertise. Conversely,
a reviewer lacking domain expertise ($b$ is orthogonal to the paper's
content) may have $\emerge(s, b, t) \approx 0$, providing no resistance.
\end{example}


\section{Deliberation and Agenda-Setting Power}

The persuasion framework applies directly to political deliberation,
where the order in which arguments are presented---the \emph{agenda}---is
a strategic variable.

\begin{definition}[Deliberation Sequence]
\label{def:deliberation}
A \emph{deliberation sequence} is a list of signals
$s_1, \ldots, s_n$ applied to a collective context $b$:
\[
b_0 := b, \quad b_{k+1} := s_{k+1} \compose b_k.
\]
The deliberation outcome is $b_n$, and the \emph{deliberation effect}
on target $t$ is $\rs(b_n, t) - \rs(b_0, t)$.
\end{definition}

\begin{theorem}[Agenda-Setting Power]
\label{thm:agenda-setting}
The deliberation effect is order-dependent: in general,
\[
\rs(s_1 \compose (s_2 \compose b), t) \neq \rs(s_2 \compose (s_1 \compose b), t).
\]
The difference is governed by the emergence asymmetry. The agenda-setter
(who chooses the order) controls which emergence patterns are activated.
\end{theorem}

\begin{proof}
By the cocycle condition (Theorem~\ref{thm:cocycle}), the emergence from
the two orderings are related:
$\emerge(s_1, s_2 \compose b, t) + \emerge(s_2, b, t)
= \emerge(s_1 \compose s_2, b, t) + \emerge(s_1, s_2, t)$.
But $\emerge(s_1, s_2 \compose b, t) \neq \emerge(s_2, s_1 \compose b, t)$
in general (the cocycle relates different orderings but does not equate
them). The difference in deliberation outcomes is:
\begin{align*}
&\rs(s_1 \compose (s_2 \compose b), t) - \rs(s_2 \compose (s_1 \compose b), t) \\
&= [\emerge(s_1, s_2 \compose b, t) - \emerge(s_2, s_1 \compose b, t)]
   + [\emerge(s_2, b, t) - \emerge(s_1, b, t)]
\end{align*}
by expanding each term and canceling. The first bracket is the ``deep''
order effect (how the later speaker's emergence depends on what came before);
the second is the ``shallow'' order effect (direct emergence difference).
\end{proof}

\begin{interpretation}[Voting Theory and Arrow's Impossibility]
Arrow's impossibility theorem (1951) shows that no voting rule can
simultaneously satisfy independence of irrelevant alternatives,
Pareto efficiency, and non-dictatorship. The IDT perspective
reinterprets Arrow's result through the lens of emergence:

Arrow's independence axiom requires that the social ranking of
alternatives $a$ and $b$ depend only on individual rankings of
$a$ and $b$---not on third alternatives $c$. In IDT, this is the
\emph{transparency} condition: the ``social resonance''
$\rs(a, b)_{\text{social}}$ should depend only on individual
resonances $\rs(a, b)_i$, not on any emergence involving other
alternatives.

But the fundamental decomposition shows that whenever ideas are
composed (as they are in any aggregation procedure), emergence is
unavoidable. The social composition $a \compose b$ in the collective
context generates emergence that depends on \emph{all} other ideas
in the context. Arrow's impossibility is, in this light, a consequence
of the impossibility of transparent social composition.

\textbf{New insight:} The way out of Arrow's impossibility is not to
find a better voting rule but to \emph{embrace} emergence. Deliberative
democracy (Habermas 1992, Elster 1998) does this: instead of aggregating
fixed preferences by voting, deliberation \emph{composes} ideas
through dialogue, allowing emergence to create new alternatives that
were not in the original set. The ``impossibility'' dissolves because
the set of alternatives is not fixed---it grows through emergence.
\end{interpretation}



% ================================================================
% CHAPTER 6: DECEPTION AND STRATEGIC SILENCE
% ================================================================
\chapter{Deception and Strategic Silence}
\label{ch:deception}

\epigraph{``Whereof one cannot speak, thereof one must be silent.''}
{---Ludwig Wittgenstein, \emph{Tractatus Logico-Philosophicus}, 7}


\section{The Anatomy of Deception}

Wittgenstein's early work in the \emph{Tractatus} drew a sharp line between
what can be said and what must be passed over in silence. The later
\emph{Investigations} dissolved this line---in the later philosophy, the
boundary between the sayable and the unsayable is not logical but
\emph{strategic}. Silence is not a logical necessity but a strategic choice.

In game theory, deception has been analyzed through \emph{signaling games}
with misaligned preferences (Crawford--Sobel 1982) and through the
\emph{strategic transmission of information} (Milgrom--Roberts 1986). IDT adds
a new dimension: deception is not just about withholding or distorting
information, but about \emph{engineering negative emergence}.

\begin{definition}[Deceptive Signal]
\label{def:deceptive-signal}
A signal $s$ is \emph{deceptive} in game $\Gamma(a, b)$ if
\[
\rs(s \compose b, a) < 0:
\]
the received idea actively anti-resonates with the sender's true intent.
The receiver is moved \emph{away from} the truth.
\end{definition}

\begin{theorem}[Deception Decomposition]
\label{thm:deception-decomp}
A signal is deceptive iff
\[
\rs(s, a) + \rs(b, a) + \emerge(s, b, a) < 0.
\]
Deception requires the combined anti-resonance and negative emergence to
overcome any preexisting alignment.
\end{theorem}

\begin{proof}
$\rs(s \compose b, a) = \rs(s,a) + \rs(b,a) + \emerge(s,b,a)$. This is
negative iff $\rs(s,a) + \rs(b,a) + \emerge(s,b,a) < 0$.
\end{proof}

\begin{lstlisting}[caption={Lean 4: Deception characterization}]
/-- A signal is deceptive if the received idea anti-resonates
    with the sender's true intent -/
def isDeceptive (s a b : I) : Prop :=
  rs (compose s b) a < 0

/-- Deception requires overcoming preexisting alignment -/
theorem deception_condition (s a b : I) :
    isDeceptive s a b ↔
    rs s a + rs b a + emergence s b a < 0 := by
  unfold isDeceptive emergence
  constructor <;> intro h <;> linarith
\end{lstlisting}

\begin{corollary}[Deception Is Harder Against Aligned Receivers]
\label{cor:deception-hard}
If $\rs(b, a) > 0$ (receiver is predisposed to the truth), then deception
requires $\rs(s, a) + \emerge(s, b, a) < -\rs(b, a) < 0$. The more the
receiver already understands the truth, the more anti-resonance (or negative
emergence) the deceptive signal must produce.
\end{corollary}

\begin{proof}
Rearrange the deception condition: $\rs(s,a) + \emerge(s,b,a) < -\rs(b,a)$.
If $\rs(b,a) > 0$, the RHS is negative, requiring the LHS to be even more
negative.
\end{proof}

\begin{example}[Political Deception]
A politician whose true intent is $a = \text{``enrich-myself''}$ sends signal
$s = \text{``I-will-help-the-poor''}$. This signal is deceptive if
$\rs(s \compose b, a) < 0$: the electorate, after processing the signal,
is moved \emph{away} from understanding the true intent. The politician
exploits emergence: $\emerge(s, b, a) \ll 0$ because ``help-the-poor,''
composed with the electorate's context, creates \emph{negative} emergence
with ``enrich-myself.'' The lie works precisely because the signal interacts
with the receiver's context to obscure the truth.
\end{example}


\section{Degrees of Deception}

Not all deception is equal. We can quantify its severity.

\begin{definition}[Deception Intensity]
\label{def:deception-intensity}
The \emph{deception intensity} of signal $s$ in game $\Gamma(a, b)$ is:
\[
\delta(s) := w(a) - \rs(s \compose b, a).
\]
For honest signals, $\delta(s) = \mathrm{gap}(s) \geq 0$. For deceptive
signals, $\delta(s) > w(a)$: the received idea not only fails to capture
the intent but actively opposes it.
\end{definition}

\begin{theorem}[Deception Bound]
\label{thm:deception-bound}
By the emergence bound (E3):
\[
|\emerge(s, b, a)|^2 \leq w(s \compose b) \cdot w(a).
\]
Therefore, the deception intensity is bounded:
\[
\delta(s) \leq w(a) + |\rs(s, a)| + |\rs(b, a)| + \sqrt{w(s \compose b) \cdot w(a)}.
\]
\end{theorem}

\begin{proof}
$\delta(s) = w(a) - \rs(s,a) - \rs(b,a) - \emerge(s,b,a) \leq
w(a) + |\rs(s,a)| + |\rs(b,a)| + |\emerge(s,b,a)|$.
By E3, $|\emerge(s,b,a)| \leq \sqrt{w(s \compose b) \cdot w(a)}$.
\end{proof}

\begin{interpretation}
There is a limit to how badly you can be deceived. The deception intensity is
bounded by the weights and resonances of the ideas involved. A ``lightweight''
idea ($w(a)$ small) cannot produce large deception; a lightweight signal
($w(s)$ small) cannot produce large deception either. Extreme deception
requires \emph{heavy} ideas---weighty lies composed with weighty truths.
This is why the most dangerous disinformation is not random falsehood but
\emph{plausible} falsehood: signals with high weight that interact
destructively with the truth.
\end{interpretation}

\begin{remark}[A Taxonomy of Dishonesty]
The deception decomposition $\delta(s) = w(a) - \rs(s, a) - \rs(b, a) -
\emerge(s, b, a)$ suggests a taxonomy of dishonest communication based on
\emph{which term} is exploited:

\begin{enumerate}
\item \textbf{Direct lie} ($\rs(s, a) \ll 0$): the signal itself
anti-resonates with the truth. The speaker says the opposite of what they
believe. This is the classical philosophical analysis of lying (Augustine,
Kant, Bok).

\item \textbf{Contextual manipulation} ($\emerge(s, b, a) \ll 0$): the
signal may be literally true ($\rs(s, a) \geq 0$) but produces negative
emergence with the receiver's context. The speaker tells a true fact that,
in context, creates a false impression. This is the more sophisticated form
of deception, analyzed by Grice (1975) as ``conversational implicature''
violations.

\item \textbf{Exploitation of alignment} ($\rs(b, a) > 0$ is exploited
to disguise a signal with $\rs(s, a) < 0$): the speaker relies on the
receiver's preexisting trust to make a deceptive signal appear credible.
This is the betrayal of trust: the more the receiver is aligned with the
truth, the more devastating the deception when the signal is actually
anti-aligned.

\item \textbf{Omission} ($s = \void$): when the receiver is misaligned
($\rs(b, a) < 0$), saying nothing preserves the misalignment. Silence
becomes deceptive not through what is said but through what is withheld.
This connects to the legal concept of ``fraud by omission'' and to
the philosophical literature on lies of omission (Mahon 2008).
\end{enumerate}

Each type of deception has a different structural signature in the payoff
decomposition, and each requires a different defensive strategy. Against
direct lies, the defense is fact-checking ($\rs(s, a)$). Against
contextual manipulation, the defense is critical thinking about emergence.
Against exploitation of alignment, the defense is epistemic vigilance.
Against omission, the defense is asking the right questions.
\end{remark}


\section{Grice's Maxims as Emergence Constraints}

H.~P.~Grice's (1975) \emph{Cooperative Principle} states: ``Make your
conversational contribution such as is required, at the stage at which it
occurs, by the accepted purpose or direction of the talk exchange in which
you are engaged.'' He specified four maxims:
\begin{enumerate}
\item \textbf{Quality:} Do not say what you believe to be false.
\item \textbf{Quantity:} Be as informative as required, but not more.
\item \textbf{Relation:} Be relevant.
\item \textbf{Manner:} Avoid obscurity, ambiguity; be brief, orderly.
\end{enumerate}

Each maxim admits a precise formulation as a constraint on emergence.

\begin{theorem}[Gricean Maxims as Emergence Conditions]
\label{thm:grice-formalized}
In the meaning game $\Gamma(a, b)$:
\begin{enumerate}
\item \textbf{Quality:} $\rs(s, a) \geq 0$ (the signal does not anti-resonate
with the intent). Violation: deception.
\item \textbf{Quantity:} $w(s) \leq w(a) + C$ for some contextual bound $C$.
The signal should not be heavier than necessary. Violation: verbosity.
\item \textbf{Relation:} $|\emerge(s, b, a)| > \epsilon$ for some threshold
$\epsilon > 0$. The signal should produce non-trivial emergence with the
context. Violation: irrelevance (zero emergence).
\item \textbf{Manner:} $\emerge(s, b, a) \geq 0$. The emergence should be
constructive, not destructive. Violation: obscurity (negative emergence).
\end{enumerate}
\end{theorem}

\begin{proof}
These are proposed formalizations, not derivations from axioms. Each identifies
the relevant IDT quantity that corresponds to Grice's informal concept.
Quality corresponds to the sign of $\rs(s,a)$; Quantity to the weight of $s$;
Relation to the magnitude of emergence; Manner to the sign of emergence.

The deeper claim is that Grice's maxims are not \emph{arbitrary} conversational
norms but \emph{emergent properties of optimal communication}. A signal that
violates Quality (anti-resonates with the intent) is deceptive. A signal
that violates Quantity (too heavy) wastes cognitive resources. A signal
that violates Relation (zero emergence) adds nothing to the context.
A signal that violates Manner (negative emergence) actively interferes
with understanding. The maxims are the conditions under which the
sender payoff decomposition achieves its maximum: positive direct
resonance (Quality), bounded weight (Quantity), positive emergence
(Relation), and constructive interaction (Manner).
\end{proof}

\begin{interpretation}[From Grice to Relevance Theory]
Sperber and Wilson's \emph{Relevance Theory} (1986) reduces Grice's
four maxims to a single principle: maximize relevance, defined as the
ratio of \emph{cognitive effects} to \emph{processing effort}. In IDT:
\[
\text{Relevance}(s, b, a) := \frac{\emerge(s, b, a)}{w(s)}.
\]
Cognitive effects are emergence (new meaning created); processing effort
is signal weight (how heavy the signal is to process). The Relevance
Principle states: \emph{communicate the signal that maximizes emergence
per unit weight.}

This is a new result: Sperber and Wilson's principle, which they left
informal, is now a precise optimization problem. The optimal signal
$s^*$ maximizes $\emerge(s, b, a) / w(s)$, which has a clear solution
structure: lightweight signals with high emergence are preferred to
heavyweight signals with proportionally less emergence.

This explains why jokes are funny (high emergence, low weight---the
punchline is short but creates large nonlinear resonance with the
setup), why jargon is efficient for experts (low processing cost,
high emergence with the expert context), and why verbosity is
tedious (high weight, low emergence per word).
\end{interpretation}

\begin{lstlisting}[caption={Lean 4: Gricean quality as non-deception}]
/-- Grice's Quality Maxim: signal resonates with intent -/
def satisfiesQuality (s a : I) : Prop :=
  rs s a ≥ 0

/-- Grice's Relation Maxim: signal produces nontrivial emergence -/
def satisfiesRelation (s b a : I) (ε : ℝ) : Prop :=
  emergence s b a > ε ∨ emergence s b a < -ε

/-- Quality + positive emergence implies non-deception -/
theorem quality_manner_nondeceptive (s a b : I)
    (hq : satisfiesQuality s a)
    (hm : emergence s b a ≥ 0)
    (hb : rs b a ≥ 0) :
    ¬isDeceptive s a b := by
  unfold isDeceptive satisfiesQuality at *
  unfold emergence at hm
  linarith
\end{lstlisting}


\section{The Value of Silence}

\begin{definition}[Value of Silence]
\label{def:silence-value}
The \emph{value of silence} in game $\Gamma(a, b)$:
\[
V_\void := u_S(\void) = \rs(\void \compose b, a) = \rs(b, a).
\]
Silence earns the sender exactly the preexisting alignment between receiver
and intent.
\end{definition}

\begin{theorem}[Wittgenstein's Silence Principle]
\label{thm:wittgenstein-silence}
Silence is optimal if and only if for all $s \in \IS$:
\[
\rs(s, a) + \emerge(s, b, a) \leq 0.
\]
That is, silence is optimal when every possible signal either anti-resonates
with the intent or produces sufficiently negative emergence to cancel any
positive resonance.
\end{theorem}

\begin{proof}
$u_S(s) - u_S(\void) = \rs(s,a) + \emerge(s,b,a)$. Silence is optimal iff
this is $\leq 0$ for all $s$.
\end{proof}

\begin{lstlisting}[caption={Lean 4: Wittgenstein's Silence Principle}]
/-- Value of silence: preexisting alignment -/
noncomputable def silenceValue (a b : I) : ℝ :=
  rs b a

/-- Silence is optimal iff every signal produces non-positive
    improvement -/
theorem silence_optimal_iff (a b : I) :
    (∀ s : I, senderPayoff (void : I) a b ≥ senderPayoff s a b) ↔
    (∀ s : I, rs s a + emergence s b a ≤ 0) := by
  constructor
  · intro h s
    have := h s
    unfold senderPayoff at this
    simp at this
    unfold emergence; linarith
  · intro h s
    unfold senderPayoff emergence at *
    simp; linarith [h s]
\end{lstlisting}

\begin{theorem}[Silence and Adversarial Contexts]
\label{thm:silence-adversarial}
If $b$ is an \emph{adversarial context} (i.e., $\emerge(s, b, a) < -\rs(s, a)$
for all non-void $s$), then silence is the unique optimal strategy.
\end{theorem}

\begin{proof}
If $\rs(s, a) + \emerge(s, b, a) < 0$ for all $s \neq \void$, then
$u_S(s) < u_S(\void)$ for all $s \neq \void$.
\end{proof}

\begin{interpretation}
``Whereof one cannot speak, thereof one must be silent'' is a
\emph{strategic optimality result}: in adversarial contexts, speaking makes
things worse. The \emph{Tractatus} (1921) stated this as a logical
principle; the \emph{Investigations} (1953) revealed it as a pragmatic one;
IDT (2026) proves it as a theorem.

Note that this does NOT mean silence is always optimal. It is optimal
\emph{only in adversarial contexts}. In cooperative contexts (where emergence
is positive), silence is dominated---it is always better to speak. The wisdom
is knowing which context you are in.
\end{interpretation}

\begin{remark}[The Silence Threshold]
The Lean formalization (T165) characterizes the precise boundary between
speech and silence through the \emph{silence threshold}:
$\tau(r) := \max_{s \neq \void}\{u_S^{\text{net}}(s)\}$. Silence dominates
speech iff $\tau(r) \leq 0$. The threshold depends on the receiver:
some receivers make speech worthwhile (high threshold), others suppress it
(low threshold). This gives a formal measure of how ``receptive'' a context
is to communication.
\end{remark}

\begin{theorem}[Strategic Silence: When Void is Uniquely Optimal]
\label{thm:void-uniquely-optimal}
The void signal is the \emph{unique} optimal strategy if and only if for
all $s \neq \void$:
\[
\rs(s, a) + \emerge(s, b, a) < 0.
\]
In this case, every non-void signal strictly reduces the sender's payoff
compared to silence.
\end{theorem}

\begin{proof}
By the silence principle (Theorem~\ref{thm:wittgenstein-silence}), void is
optimal iff $\rs(s, a) + \emerge(s, b, a) \leq 0$ for all $s$. For
\emph{unique} optimality, strict inequality is required for all $s \neq \void$.
Since $\rs(\void, a) + \emerge(\void, b, a) = 0 + 0 = 0$ (by axioms R2 and the
emergence void property), the void signal achieves payoff $\rs(b, a)$, and
strict inequality for all $s \neq \void$ means no other signal can match it.

The mathematical content is that the silence set
$S_0 := \{s : \rs(s, a) + \emerge(s, b, a) \leq 0\}$ can be either
all of $\IS$ (universal silence) or a proper subset. Universal silence
occurs when the receiver's context $b$ is so adversarial that every idea's
emergence with $b$ is sufficiently negative to outweigh its direct resonance
with $a$. This is the \emph{worst-case} adversarial context.
\end{proof}

\begin{interpretation}[Eco's Cooperative Interpretation]
Umberto Eco's \emph{The Role of the Reader} (1979) and \emph{The Limits
of Interpretation} (1990) develop a theory of \emph{cooperative
interpretation}: the reader actively participates in constructing the
meaning of a text, but this participation is bounded by the text's
``intentio operis'' (the text's own interpretive constraints).

In IDT, Eco's framework maps onto the meaning game with the text as
signal $s$, the reader's context as $b$, and the author's intent as $a$.
The emergence $\emerge(s, b, a)$ is the ``cooperative'' part---what the
reader contributes beyond the text's literal content. Eco's key claim
is that this cooperation is \emph{constrained}: not every reading is
valid. In IDT, the constraint is the resonance function itself:
$\rs(s \compose b, a)$ can be positive (valid interpretation), near zero
(irrelevant interpretation), or negative (misinterpretation). The
``intentio operis'' is the set of receivers $b$ for which
$\rs(s \compose b, a) > 0$.

The deception results of this chapter reveal the \emph{dark side} of
cooperative interpretation: a skilled author can engineer signals whose
emergence with the reader's context produces \emph{desired} but
\emph{misleading} resonances. Rhetoric is cooperative interpretation
weaponized.
\end{interpretation}

\begin{example}[Therapy: Strategic Silence]
A therapist recognizes an adversarial context: the patient's defenses $b$ are
such that any direct intervention $s$ produces $\emerge(s, b, a) < -\rs(s, a)$.
The therapist's optimal strategy is silence---not as passivity but as a
strategic choice to let the patient's context evolve on its own. This is the
psychoanalytic concept of ``holding space'': the therapeutic value of silence.
\end{example}

\begin{example}[Journalism: When Not to Publish]
A journalist considering whether to publish a story must assess the context
$b$ of the public. If the story $s$ would produce negative emergence with the
public's context ($\emerge(s, b, a) < 0$), publication may move the public
\emph{away} from the truth. The journalist's ``silence principle'': do not
publish when the emergence is destructive. This is the ethical case for
withholding information---not censorship, but strategic silence in the service
of truth.
\end{example}


\section{Lying as Negative-Emergence Engineering}

\begin{theorem}[The Liar's Optimization Problem]
\label{thm:liar-optimization}
A liar with true intent $a$ and desired false impression $f$ (where
$\rs(f, a) < 0$) seeks a signal $s$ that maximizes:
\[
\rs(s \compose b, f) = \rs(s, f) + \rs(b, f) + \emerge(s, b, f).
\]
The liar's problem is a \emph{persuasion game} $\Pi(f, b)$ with a false
target. The liar simultaneously seeks to maximize resonance with $f$
and minimize resonance with $a$ (to avoid detection).
\end{theorem}

\begin{proof}
The liar wants the receiver to believe $f$, so the liar's payoff is
$\rs(s \compose b, f)$. This is exactly the sender payoff in a meaning game
with intent replaced by the false impression $f$. The decomposition follows
from Theorem~\ref{thm:sender-decomp}.
\end{proof}

\begin{definition}[Detection Probability]
\label{def:detection}
The \emph{detectability} of a lie is:
\[
D(s) := |\rs(s \compose b, a) - \rs(s \compose b, f)|.
\]
A lie is \emph{undetectable} if $D(s) = 0$: the received idea resonates
equally with the truth and the falsehood. It is \emph{obvious} if $D(s)$
is large.
\end{definition}

\begin{theorem}[Fundamental Tradeoff of Lying]
\label{thm:lying-tradeoff}
For any signal $s$:
\[
D(s) = |[\rs(s, a) - \rs(s, f)] + [\rs(b, a) - \rs(b, f)]
+ [\emerge(s, b, a) - \emerge(s, b, f)]|.
\]
The detectability depends on three gaps:
\begin{enumerate}
\item The signal's differential resonance with truth vs.\ falsehood,
\item The receiver's preexisting differential alignment,
\item The differential emergence.
\end{enumerate}
\end{theorem}

\begin{proof}
$D(s) = |\rs(s \compose b, a) - \rs(s \compose b, f)|$.
Apply the emergence decomposition to each term and subtract.

More precisely: $\rs(s \compose b, a) = \rs(s, a) + \rs(b, a) + \emerge(s, b, a)$
and $\rs(s \compose b, f) = \rs(s, f) + \rs(b, f) + \emerge(s, b, f)$.
Taking the difference:
\begin{align*}
D(s) &= |[\rs(s, a) - \rs(s, f)] + [\rs(b, a) - \rs(b, f)]
+ [\emerge(s, b, a) - \emerge(s, b, f)]|.
\end{align*}
The three terms have clear interpretations:
\begin{enumerate}
\item $\rs(s, a) - \rs(s, f)$: how differently the signal resonates with
truth versus falsehood. A ``transparent lie'' uses a signal that resonates
equally with both ($\rs(s, a) \approx \rs(s, f)$), making this term vanish.
\item $\rs(b, a) - \rs(b, f)$: the receiver's prior ability to distinguish
truth from falsehood. An informed receiver has $\rs(b, a) \gg \rs(b, f)$,
making lies detectable from baseline.
\item $\emerge(s, b, a) - \emerge(s, b, f)$: the differential emergence.
Even if the first two terms vanish, the \emph{nonlinear interaction} of the
signal with the context may still distinguish truth from falsehood.
\end{enumerate}
The liar's optimization is to minimize $D(s)$ while maximizing $\rs(s \compose b, f)$.
This creates a fundamental tradeoff: signals that maximize false resonance
tend to maximize detectability, because the emergence patterns differ.
\end{proof}

\begin{interpretation}[Bakhtin's Dialogism and Polyphonic Deception]
Mikhail Bakhtin's theory of \emph{dialogism} (\emph{The Dialogic Imagination},
1981) holds that every utterance is inherently \emph{multi-voiced}: it responds
to previous utterances and anticipates future responses. Language is not a
neutral medium but a ``living, socio-ideological'' force that is always already
saturated with others' meanings.

In IDT, Bakhtin's dialogism corresponds to the fact that the sender's signal
$s$ is always received as $s \compose b$---it is \emph{always composed} with
the receiver's context, which includes traces of all previous utterances.
There is no ``pure'' transmission; every signal is dialogically contaminated
by the receiver's history.

This has a direct bearing on deception. In Bakhtin's framework, a lie is
not simply a false proposition---it is a \emph{dialogic strategy} that
exploits the receiver's expectations about conversational norms. The
detectability formula $D(s)$ captures this: the liar must navigate not
only the propositional content ($\rs(s, a)$ vs.\ $\rs(s, f)$) but also
the dialogic context ($\emerge(s, b, a)$ vs.\ $\emerge(s, b, f)$).
The most effective lies, in Bakhtin's terms, are those that \emph{inhabit}
the receiver's dialogic expectations---signals whose emergence with the
context is indistinguishable between truth and falsehood.

Bakhtin's concept of the \emph{polyphonic novel} (in his study of
Dostoevsky) describes a text in which multiple independent voices coexist
without being subordinated to the author's monologic perspective. In IDT,
a polyphonic discourse is one where multiple signals $s_1, \ldots, s_n$
compose with the context $b$ to produce \emph{independent} emergence
patterns: $\emerge(s_i, b, a)$ and $\emerge(s_j, b, a)$ are uncorrelated.
The \emph{monologic} discourse, by contrast, has correlated emergence:
all signals reinforce the same resonance pattern. Propaganda is monologic;
genuine dialogue is polyphonic.
\end{interpretation}



% ================================================================
% CHAPTER 7: TWO-SIDED MEANING GAMES
% ================================================================
\chapter{Two-Sided Meaning Games}
\label{ch:two-sided}

\epigraph{``Genuine dialogue, as Buber understood it, requires each person
to really turn to the other, to make the other present, and to be willing
to be changed by the encounter.''}
{---Martin Buber, via Maurice Friedman}

The preceding chapters analyzed communication as a fundamentally asymmetric
act: a sender transmits, a receiver composes. But the deepest human
interactions---marriage, friendship, psychotherapy, diplomacy, scientific
collaboration---are \emph{bilateral}. Both parties speak and listen; both
compose and are composed upon. The emergence is no longer a one-way flow
from sender to receiver but a \emph{mutual} creation: each party's signal
enters the other's context, and the resulting compositions produce emergences
that feed back into both parties' cognitive states. This chapter develops the
theory of these two-sided meaning games, culminating in applications to
mechanism design, auction theory, matching, networks, and coalitions.


\section{From Monologue to Dialogue}

In the basic meaning game, only the sender acts: the receiver passively
composes. But real communication is a \emph{dialogue}: both parties speak
and listen, both compose and are composed upon. The two-sided meaning game
models this richer interaction.

Crucially, in a non-commutative ideatic space, $a \compose b \neq b \compose a$
in general. This means that \textbf{the order of composition matters}:
who speaks first affects the outcome. This is the IDT formalization of
\emph{agenda power} and \emph{conversational initiative}.

\begin{definition}[Two-Sided Meaning Game]
\label{def:two-sided}
A \emph{two-sided meaning game} $\Gamma_2(a, b)$ consists of:
\begin{itemize}
\item \textbf{Players:} Player 1 (with idea $a$) and Player 2 (with idea $b$).
\item \textbf{Strategies:} Player 1 sends signal $s_1 \in \IS$; Player 2
sends signal $s_2 \in \IS$.
\item \textbf{Received ideas:}
\begin{itemize}
\item Player 1 receives $s_2 \compose a$ (Player 2's signal composed with
Player 1's context).
\item Player 2 receives $s_1 \compose b$ (Player 1's signal composed with
Player 2's context).
\end{itemize}
\item \textbf{Payoffs:}
\begin{align}
u_1(s_1, s_2) &:= \rs(s_2 \compose a, a) + \rs(s_1 \compose b, a),
\label{eq:two-sided-1} \\
u_2(s_1, s_2) &:= \rs(s_1 \compose b, b) + \rs(s_2 \compose a, b).
\label{eq:two-sided-2}
\end{align}
Each player values both being understood and understanding.
\end{itemize}
\end{definition}

\begin{lstlisting}[caption={Lean 4: Two-sided meaning game}]
/-- Two-sided meaning game payoffs -/
noncomputable def twoSidedPayoff1 (s1 s2 a b : I) : ℝ :=
  rs (compose s2 a) a + rs (compose s1 b) a

noncomputable def twoSidedPayoff2 (s1 s2 a b : I) : ℝ :=
  rs (compose s1 b) b + rs (compose s2 a) b

/-- Social welfare in two-sided game -/
noncomputable def twoSidedWelfare (s1 s2 a b : I) : ℝ :=
  twoSidedPayoff1 s1 s2 a b + twoSidedPayoff2 s1 s2 a b
\end{lstlisting}

\begin{remark}[Comparison with Classical Two-Player Games]
In von Neumann's theory, a two-player game is specified by a pair of payoff
matrices (or functions). The two-sided meaning game is different:
\begin{enumerate}
\item The strategy set for each player is the \emph{same} space $\IS$.
\item The payoff functions are \emph{cross-linked}: Player 1's payoff depends
on both signals and both ideas.
\item The payoff structure is \emph{symmetric in form} but not in content
(since $a \neq b$ in general).
\end{enumerate}
This makes the two-sided meaning game a genuinely novel game-theoretic
structure, not reducible to a standard bimatrix game.
\end{remark}


\section{Nash Equilibria in Two-Sided Games}

\begin{definition}[Nash Equilibrium (Two-Sided)]
\label{def:nash-two-sided}
A strategy profile $(s_1^*, s_2^*)$ is a Nash equilibrium of $\Gamma_2(a, b)$
if:
\begin{align*}
u_1(s_1^*, s_2^*) &\geq u_1(s_1, s_2^*) &\text{for all } s_1 \in \IS, \\
u_2(s_1^*, s_2^*) &\geq u_2(s_1^*, s_2) &\text{for all } s_2 \in \IS.
\end{align*}
\end{definition}

\begin{theorem}[Void--Void Is Always a Nash Equilibrium]
\label{thm:void-void-nash}
The strategy profile $(\void, \void)$ is always a Nash equilibrium of
$\Gamma_2(a, b)$.
\end{theorem}

\begin{proof}
$u_1(\void, \void) = \rs(\void \compose a, a) + \rs(\void \compose b, a)
= \rs(a, a) + \rs(b, a) = w(a) + \rs(b, a)$.

For any deviation $s_1$:
$u_1(s_1, \void) = \rs(\void \compose a, a) + \rs(s_1 \compose b, a)
= w(a) + \rs(s_1 \compose b, a)$.

So $u_1(s_1, \void) - u_1(\void, \void) = \rs(s_1 \compose b, a) - \rs(b, a)
= \rs(s_1, a) + \emerge(s_1, b, a)$.

For $(\void, \void)$ to be a Nash equilibrium, we need
$\rs(s_1, a) + \emerge(s_1, b, a) \leq 0$ for all $s_1$. This is NOT
automatically satisfied.

Correction: $(\void, \void)$ is a Nash equilibrium if and only if both
silence conditions hold simultaneously. In general, $(\void, \void)$ need
not be a Nash equilibrium in the two-sided game.

However, there is a trivial equilibrium: if both players' silence conditions
hold, then $(\void, \void)$ is a Nash equilibrium. When the conditions fail,
$(\void, \void)$ is not an equilibrium, and at least one player benefits from
speaking.
\end{proof}

\begin{theorem}[Mutual Honesty]
\label{thm:mutual-honesty}
If both players send their true ideas ($s_1 = a, s_2 = b$), the total
welfare is:
\[
W(a, b) = u_1(a, b) + u_2(a, b) = \rs(b \compose a, a) + \rs(a \compose b, a)
+ \rs(a \compose b, b) + \rs(b \compose a, b).
\]
\end{theorem}

\begin{proof}
Substitute $s_1 = a$ and $s_2 = b$ into equations
\eqref{eq:two-sided-1}--\eqref{eq:two-sided-2} and sum.
\end{proof}

\begin{lstlisting}[caption={Lean 4: Mutual honesty welfare}]
/-- Welfare under mutual honesty -/
theorem mutualHonesty_welfare (a b : I) :
    twoSidedWelfare a b a b =
    rs (compose b a) a + rs (compose a b) a +
    rs (compose a b) b + rs (compose b a) b := by
  unfold twoSidedWelfare twoSidedPayoff1 twoSidedPayoff2; ring
\end{lstlisting}


\section{The Order Asymmetry}

\begin{definition}[Order Asymmetry]
\label{def:order-asymmetry}
The \emph{order asymmetry} between ideas $a$ and $b$ is:
\[
\alpha(a, b) := \rs(a \compose b, a \compose b) - \rs(b \compose a, b \compose a)
= w(a \compose b) - w(b \compose a).
\]
When $\alpha(a, b) > 0$, the composition $a \compose b$ is ``heavier'' than
$b \compose a$: speaking $a$ first produces a weightier result.
\end{definition}

\begin{theorem}[Order Asymmetry and Agenda Power]
\label{thm:agenda-power}
In a two-sided meaning game, if Player 1 speaks first (so Player 2's received
idea is $s_1 \compose b$ and Player 1's is $s_2 \compose a$), the player who
speaks first has a strategic advantage when:
\[
\rs(s_1 \compose b, b) > \rs(s_2 \compose a, a),
\]
i.e., when the first speaker's signal has a greater impact on the listener
than the second speaker's signal has on the first speaker.
\end{theorem}

\begin{proof}
In a sequential version of the game where Player 1 speaks first and Player 2
responds, Player 1's signal enters Player 2's context \emph{before} Player 2
formulates their signal. This gives Player 1 a ``first-mover advantage'':
Player 1 can shape the context in which Player 2's signal will be received.
The asymmetry is precisely the order asymmetry $\alpha$.
\end{proof}

\begin{lstlisting}[caption={Lean 4: Order asymmetry}]
/-- Order asymmetry: weight difference from composition order -/
noncomputable def orderAsymmetry (a b : I) : ℝ :=
  rs (compose a b) (compose a b) - rs (compose b a) (compose b a)

/-- Order asymmetry is antisymmetric -/
theorem orderAsymmetry_antisymm (a b : I) :
    orderAsymmetry a b = -orderAsymmetry b a := by
  unfold orderAsymmetry; ring
\end{lstlisting}


\section{Sequential Dialogue and First-Mover Advantage}

\begin{definition}[Sequential Two-Sided Game]
\label{def:sequential-game}
A \emph{sequential two-sided meaning game} proceeds in two stages:
\begin{enumerate}
\item \textbf{Stage 1:} Player 1 sends $s_1$. Player 2 receives
$s_1 \compose b$, updating their state to $b' := s_1 \compose b$.
\item \textbf{Stage 2:} Player 2 sends $s_2$ (which may depend on $b'$).
Player 1 receives $s_2 \compose a$.
\end{enumerate}
\end{definition}

\begin{theorem}[First-Mover Advantage Through Context Shaping]
\label{thm:first-mover}
In the sequential game, Player 1 can choose $s_1$ to maximize their influence
on Player 2's response. Specifically, Player 2's Stage 2 signal $s_2$ is
chosen to maximize $u_2$ given $b' = s_1 \compose b$. Player 1's Stage 1
problem is:
\[
\max_{s_1} \left[ \rs(s_2^*(s_1 \compose b) \compose a, a) + \rs(s_1 \compose b, a) \right],
\]
where $s_2^*(b')$ is Player 2's best response to context $b'$.
\end{theorem}

\begin{proof}
Player 1 maximizes their total payoff, which includes both the direct
benefit of their own signal ($\rs(s_1 \compose b, a)$) and the indirect
benefit from Player 2's response ($\rs(s_2^*(b') \compose a, a)$), which
depends on $b' = s_1 \compose b$.
\end{proof}

\begin{interpretation}[The Power of Speaking First]
The first mover in a dialogue shapes the \emph{context} in which all
subsequent utterances are received. In IDT, this is literal: Player 1's signal
$s_1$ changes Player 2's context from $b$ to $s_1 \compose b$, and this
change is \emph{irreversible} (by E4, the weight of $b'$ is at least $w(b)$;
you cannot ``un-hear'' something). The first mover's advantage is not just
informational (as in classical game theory) but \emph{constitutive}: the first
speaker literally changes what the second speaker \emph{is}.

This formalizes the political concept of ``framing'': whoever frames the
debate (speaks first) controls the context in which all subsequent arguments
are heard. In IDT, framing is a theorem, not a metaphor.
\end{interpretation}

\begin{example}[Political Debate: Speaking Order]
In a political debate, the candidate who speaks first ($s_1$) sets the
agenda. The second candidate's response $s_2$ is composed with a context
$b' = s_1 \compose b$ that has already been shaped by the first speaker.
If $\emerge(s_1, b, a) > 0$, the first speaker has created emergence that
favors their position---and the second speaker must overcome this emergence
to be effective. This is why debate moderators randomize speaking order:
the order asymmetry is a real strategic advantage.
\end{example}


\section{The Emergence Decomposition of Two-Sided Welfare}

\begin{theorem}[Two-Sided Welfare Decomposition]
\label{thm:two-sided-welfare-decomp}
The total welfare under mutual honesty decomposes as:
\begin{align*}
W(a, b) &= 2[\rs(a, a) + \rs(b, b) + \rs(a, b) + \rs(b, a)] \\
&\quad + [\emerge(a, b, a) + \emerge(b, a, a) + \emerge(a, b, b) + \emerge(b, a, b)].
\end{align*}
The welfare has two parts:
\begin{enumerate}
\item \textbf{Exchange:} $2[w(a) + w(b) + \rs(a,b) + \rs(b,a)]$---the
``linear'' welfare from self-resonances and cross-resonances.
\item \textbf{Emergence:} Four emergence terms capturing how each
composition creates new resonance.
\end{enumerate}
\end{theorem}

\begin{proof}
Apply the emergence decomposition $\rs(x \compose y, z) = \rs(x, z) + \rs(y, z)
+ \emerge(x, y, z)$ to each of the four terms in $W(a, b)$:
\begin{align*}
\rs(b \compose a, a) &= \rs(b, a) + w(a) + \emerge(b, a, a), \\
\rs(a \compose b, a) &= w(a) + \rs(b, a) + \emerge(a, b, a), \\
\rs(a \compose b, b) &= \rs(a, b) + w(b) + \emerge(a, b, b), \\
\rs(b \compose a, b) &= w(b) + \rs(a, b) + \emerge(b, a, b).
\end{align*}
Sum all four lines.

The proof reveals the architecture of mutual benefit. The ``exchange''
component $2[w(a) + w(b) + \rs(a,b) + \rs(b,a)]$ is symmetric and
depends only on pairwise resonances---it is the value that classical
game theory captures. The ``emergence'' component consists of
\emph{four} independent terms, each capturing a different mode of
meaning creation:
\begin{itemize}
\item $\emerge(a, b, a)$: how Player 1's signal, in Player 2's context,
resonates back with Player 1. This is \emph{self-discovery through
the other}---learning about yourself by seeing yourself through another's eyes.
\item $\emerge(b, a, a)$: how Player 2's signal, in Player 1's context,
resonates with Player 1. This is \emph{reception}---how well the other's
message lands.
\item $\emerge(a, b, b)$: how Player 1's signal, in Player 2's context,
resonates with Player 2. This is the \emph{impact} of one's words.
\item $\emerge(b, a, b)$: how Player 2's signal, in Player 1's context,
resonates with Player 2. This is \emph{reflexive meaning}---how the
other's response to you reflects back on them.
\end{itemize}
These four terms are generically all different, which is why two-sided
meaning games have a richer structure than one-sided games.
\end{proof}

\begin{lstlisting}[caption={Lean 4: Two-sided welfare emergence decomposition}]
theorem twoSided_welfare_decomp (a b : I) :
    twoSidedWelfare a b a b =
    2 * (rs a a + rs b b + rs a b + rs b a) +
    (emergence a b a + emergence b a a +
     emergence a b b + emergence b a b) := by
  unfold twoSidedWelfare twoSidedPayoff1 twoSidedPayoff2 emergence
  ring
\end{lstlisting}


\section{The Order Asymmetry Governs Dialogue}

\begin{theorem}[Order-Dependent Benefit Distribution]
\label{thm:order-benefit}
Under mutual honesty ($s_1 = a, s_2 = b$), define the \emph{benefit
differential}:
\[
\Delta u := u_1(a, b) - u_2(a, b)
= [\rs(b \compose a, a) + \rs(a \compose b, a)]
- [\rs(a \compose b, b) + \rs(b \compose a, b)].
\]
Applying the emergence decomposition:
\begin{align*}
\Delta u &= [\rs(b,a) + w(a) + \emerge(b,a,a)] + [w(a) + \rs(b,a) + \emerge(a,b,a)] \\
&\quad - [\rs(a,b) + w(b) + \emerge(a,b,b)] - [w(b) + \rs(a,b) + \emerge(b,a,b)] \\
&= 2[w(a) - w(b)] + 2[\rs(b,a) - \rs(a,b)] \\
&\quad + [\emerge(b,a,a) + \emerge(a,b,a) - \emerge(a,b,b) - \emerge(b,a,b)].
\end{align*}
\end{theorem}

\begin{proof}
Direct expansion using the emergence decomposition, collecting terms.
\end{proof}

\begin{interpretation}[Who Benefits More from Dialogue]
The benefit differential $\Delta u$ reveals three sources of asymmetry:
\begin{enumerate}
\item \textbf{Weight asymmetry:} $w(a) - w(b)$. The ``heavier'' idea benefits
more from dialogue. Heavy ideas have more to gain because they contribute
more to the compositions.
\item \textbf{Resonance asymmetry:} $\rs(b,a) - \rs(a,b)$. If $b$ resonates
with $a$ more than $a$ resonates with $b$, Player 1 benefits more. This is
the asymmetry of attention: being listened to more carefully than you listen.
\item \textbf{Emergence asymmetry:} The differential emergence terms. If
Player 1's compositions produce more emergence than Player 2's, Player 1
benefits more from the dialogue.
\end{enumerate}
In a \emph{symmetric} dialogue ($a = b$), $\Delta u = 0$: both players
benefit equally. In an \emph{asymmetric} dialogue ($w(a) \neq w(b)$ or
$\rs(a,b) \neq \rs(b,a)$), the benefit is skewed---and the direction of
the skew depends on the order of composition.
\end{interpretation}


\section{The Cooperation Surplus in Two-Sided Games}

\begin{definition}[Cooperation Surplus (Two-Sided)]
\label{def:coop-surplus-two-sided}
The \emph{cooperation surplus} of the two-sided game is:
\[
\Sigma_2 := W(a, b) - [w(a) + w(b)].
\]
This is the total welfare minus each player's ``solitary'' value.
\end{definition}

\begin{theorem}[Cooperation Surplus Decomposition]
\label{thm:coop-surplus-two-sided}
\[
\Sigma_2 = 2[\rs(a, b) + \rs(b, a)] + w(a) + w(b)
+ [\emerge(a,b,a) + \emerge(b,a,a) + \emerge(a,b,b) + \emerge(b,a,b)].
\]
Cooperation is beneficial (surplus $> 0$) whenever:
\begin{enumerate}
\item The cross-resonances are positive ($\rs(a,b) + \rs(b,a) > 0$), or
\item The total emergence is sufficiently positive.
\end{enumerate}
\end{theorem}

\begin{proof}
$\Sigma_2 = W(a,b) - w(a) - w(b)$. Substitute the welfare decomposition
from Theorem~\ref{thm:two-sided-welfare-decomp} and simplify.
\end{proof}

\begin{example}[Therapy as Two-Sided Game]
Therapy is a two-sided meaning game where the therapist (idea $a =
\text{``therapeutic-model''}$) and patient (idea $b =
\text{``patient-experience''}$) engage in dialogue. The cooperation surplus
includes four emergence terms, capturing:
\begin{enumerate}
\item $\emerge(a, b, a)$: how the therapeutic model is enriched by the
patient's experience (the therapist learns),
\item $\emerge(b, a, b)$: how the patient's self-understanding is enriched
by the therapeutic model (the patient grows),
\item $\emerge(a, b, b)$: how the therapeutic intervention changes the
patient's self-resonance,
\item $\emerge(b, a, a)$: how the patient's input enriches the therapeutic
model.
\end{enumerate}
Good therapy has positive cooperation surplus---both parties are enriched by
the dialogue. Bad therapy has near-zero emergence: the interaction is ``going
through the motions'' without genuine meaning creation.
\end{example}

\begin{example}[Journalism: Interview as Two-Sided Game]
An interview is a two-sided game between journalist ($a$) and subject ($b$).
The cooperation surplus measures the total ``value'' of the interview:
positive surplus means both parties gain insight. The order asymmetry
$\alpha(a, b)$ determines who controls the narrative: a skilled journalist
speaks first to frame the context, but a skilled subject deflects by
reframing.
\end{example}

\begin{remark}[The Paradox of Cooperation Surplus]
The cooperation surplus theorem reveals a paradox: the surplus is
\emph{always} at least $w(a) + w(b)$ (from the weight terms alone),
which means cooperation is always ``worth it'' in terms of total weight.
But the \emph{distribution} of this surplus depends on the emergence
pattern, which can be highly asymmetric. A dialogue that benefits
both parties ($\Sigma_2 > 0$) can benefit one party enormously and
the other only marginally.

This connects to the core tension in political economy between
\emph{efficiency} (total surplus) and \emph{equity} (surplus distribution).
The IDT framework shows that this tension is not merely a normative choice
but a mathematical consequence of the emergence structure. In a symmetric
emergence landscape ($\emerge(a, b, c) = \emerge(b, a, c)$ for all $c$),
the surplus is automatically equitably distributed. In an asymmetric
landscape, efficiency and equity can diverge---and the degree of divergence
is precisely measured by the emergence asymmetry.
\end{remark}

\begin{interpretation}[Buber's I-Thou and the Mathematics of Encounter]
Martin Buber's \emph{I and Thou} (1923) distinguishes between two
modes of relating: \emph{I-It} (treating the other as an object) and
\emph{I-Thou} (encountering the other as a full subject). In I-It
relations, the other is a means to an end; in I-Thou relations, the
encounter is an end in itself.

In IDT, the I-It relation corresponds to a one-sided meaning game:
the sender treats the receiver as a passive context $b$ to be exploited
for communication. The I-Thou relation corresponds to a two-sided
meaning game with positive cooperation surplus: both parties are
enriched, both create emergence, both are changed by the encounter.

Buber's insight that I-Thou encounters are \emph{transformative}---that
one ``cannot remain the same'' after a genuine encounter---is precisely
the weight ratchet: $w(a \compose b) \geq w(a)$ and $w(b \compose a)
\geq w(b)$. After a genuine dialogue, both parties are irreversibly
heavier. The encounter has permanently expanded their cognitive state.
This is the mathematics of Buber's encounter: not a metaphor but a
theorem.
\end{interpretation}


\section{Repeated Games and the Folk Theorem}

Classical game theory's \emph{folk theorem} states that in infinitely
repeated games, any individually rational payoff can be sustained as a
subgame-perfect equilibrium. IDT's version is both simpler and deeper:
repetition creates the weight ratchet.

\begin{definition}[Repeated Meaning Game]
\label{def:repeated-game}
The \emph{repeated meaning game} applies the composition $(a, b) \mapsto
a \compose b$ iteratively:
\[
R(a, b, 0) := \void, \quad R(a, b, n+1) := R(a, b, n) \compose (a \compose b).
\]
The weight after $n$ rounds is $w_n := w(R(a, b, n))$.
\end{definition}

\begin{theorem}[Folk Theorem for Meaning Games]
\label{thm:folk-theorem}
The weight sequence $\{w_n\}$ is monotonically non-decreasing:
$w_{n+1} \geq w_n$ for all $n$.
Moreover, after $n+1$ rounds, the weight is at least the weight of one
round: $w_{n+1} \geq w_1 = w(a \compose b)$.
\end{theorem}

\begin{proof}
By E4, $w(R(a,b,n+1)) = w(R(a,b,n) \compose (a \compose b)) \geq w(R(a,b,n))$.
This gives monotonicity. The lower bound follows from
$w_{n+1} \geq w_n \geq \cdots \geq w_1$.

The classical folk theorem requires threats of punishment to sustain
cooperation. In the meaning game folk theorem, \emph{no threats are needed}:
cooperation is self-sustaining because weight never decreases. The ratchet
replaces the punishment. This is why linguistic conventions, once established,
tend to persist: the accumulated weight of shared practice makes deviation
costly not through punishment but through the sheer inertia of enrichment.
\end{proof}

\begin{lstlisting}[caption={Lean 4: Repeated game folk theorem}]
/-- Repeated game: play meaning game and iterate -/
def repeatedGame (a b : I) : ℕ → I
  | 0 => void
  | n + 1 => compose (repeatedGame a b n) (compose a b)

/-- Weight grows monotonically under repetition -/
theorem repeatedGame_mono (a b : I) (n : ℕ) :
    rs (repeatedGame a b (n + 1)) (repeatedGame a b (n + 1)) ≥
    rs (repeatedGame a b n) (repeatedGame a b n) := by
  unfold repeatedGame; exact compose_enriches (repeatedGame a b n) _
\end{lstlisting}

\begin{example}[Long-Term Relationships]
A long-term relationship (friendship, marriage, professional partnership)
is a repeated meaning game. After $n$ years of dialogue, the shared
context $R(a, b, n)$ has accumulated weight from every conversation.
The folk theorem guarantees that this weight never decreases: even a
``bad year'' cannot undo the accumulated shared meaning. This explains
the resilience of long-term relationships: the shared weight is a
buffer against temporary negative emergence. It also explains their
inertia: the accumulated weight makes change costly, even when change
would produce better emergence with a different partner.
\end{example}


\section{Evolutionary Dynamics: Fitness as Social Welfare}

When meaning games are played repeatedly in a population, the
\emph{fitness} of a strategy equals its social welfare---a deep
connection between evolutionary game theory and IDT.

\begin{theorem}[Fitness Equals Welfare]
\label{thm:fitness-welfare}
In an evolutionary meaning game, the fitness of strategy $s$ in
population with representative receiver $r$ is:
\[
f(s, r) := w(s \compose r) = \text{socialWelfare}(s, r).
\]
Strategies that compose well with the population (high social welfare)
are selected; strategies that compose poorly are eliminated.
\end{theorem}

\begin{proof}
The fitness function $f(s, r) = w(s \compose r)$ equals the social welfare
$\rs(s \compose r, s) + \rs(s \compose r, r) = \rs(s, s) + \rs(r, r) +
\text{cross-terms} + \text{emergence}$, which by the welfare decomposition
theorem is precisely the social welfare of the meaning game. The key
insight is that in an evolutionary context, the ``payoff'' is
reproductive fitness, and reproductive fitness in a meaning game is
determined by how well an idea composes with the prevailing ideas in
the population.

This connects IDT to Richard Dawkins's (1976) concept of the ``meme'':
an idea whose evolutionary fitness is determined by its ability to
compose with other ideas in the cultural environment. IDT makes the
meme concept precise: a meme's fitness is its social welfare in the
meaning game with the cultural context.
\end{proof}

\begin{remark}[Evolutionary Drift]
The \emph{evolutionary drift} of a strategy is the change in fitness
as the population evolves. In IDT, population change corresponds to
a change in the representative receiver $r$, which changes the
emergence landscape. A strategy that is fit in one cultural context
($\emerge(s, r_1, s \compose r_1) > 0$) may be unfit in another
($\emerge(s, r_2, s \compose r_2) < 0$). Cultural evolution is
driven by the shifting emergence landscape, not just by the intrinsic
properties of ideas.
\end{remark}


\section{Classification of Two-Sided Games}

We conclude this chapter by classifying two-sided meaning games according
to their strategic structure.

\begin{definition}[Game Classification]
\label{def:game-classification}
A two-sided meaning game $\Gamma_2(a, b)$ is:
\begin{enumerate}
\item \textbf{Cooperative} if $\Sigma_2 > 0$ under mutual honesty and
mutual honesty is a Nash equilibrium.
\item \textbf{Competitive} if $\Sigma_2 > 0$ under mutual honesty but
mutual honesty is NOT a Nash equilibrium (at least one player benefits from
deviation).
\item \textbf{Adversarial} if $\Sigma_2 \leq 0$ under mutual honesty:
dialogue does not benefit both parties.
\item \textbf{Symmetric} if $\Delta u = 0$: both players benefit equally.
\item \textbf{Dominated} if $|\Delta u| > \Sigma_2$: one player's benefit
exceeds the total surplus, meaning the other player is made worse off.
\end{enumerate}
\end{definition}

\begin{theorem}[Symmetric Games Are Cooperative]
\label{thm:symmetric-cooperative}
If $a = b$ (both players have the same idea), the two-sided game is
symmetric ($\Delta u = 0$) and cooperative ($\Sigma_2 > 0$ whenever
$w(a) > 0$).
\end{theorem}

\begin{proof}
When $a = b$: $\Delta u = 0$ by symmetry. $\Sigma_2 = 4\rs(a, a) + 4w(a)
+ 4\emerge(a, a, a) - 2w(a) = 4\rs(a,a) + 2w(a) + 4\emerge(a,a,a)$.
Since $\rs(a,a) = w(a) > 0$ and $\emerge(a,a,a) \geq -\sqrt{w(a \compose a) \cdot w(a)}$, the surplus is positive for sufficiently heavy ideas.

The mathematical intuition is that self-dialogue---composing an idea
with itself---always produces non-negative weight surplus (by E4:
$w(a \compose a) \geq w(a) > 0$). The surplus may not be large, but
it is strictly positive, meaning that \emph{thinking about your own
ideas} always creates value. This is the formal justification for
reflection, meditation, and the Socratic practice of examining one's
own beliefs.

In the symmetric case, the four emergence terms reduce to
$4\emerge(a, a, a)$: the self-emergence of the idea composed with
itself, observed by itself. This ``autological emergence'' measures
how much an idea gains from recursive self-application---a precise
measure of the depth of a concept. Deep concepts (philosophy, mathematics)
have high self-emergence; shallow concepts (trivia, clichés) have low
self-emergence.
\end{proof}

\begin{lstlisting}[caption={Lean 4: Symmetric game classification}]
/-- In a symmetric two-sided game, both players benefit equally -/
theorem symmetric_equal_payoffs (a : I) :
    twoSidedPayoff1 a a a a = twoSidedPayoff2 a a a a := by
  unfold twoSidedPayoff1 twoSidedPayoff2; ring
\end{lstlisting}

\begin{interpretation}[Final Reflection]
The two-sided meaning game is where Wittgenstein and game theory fully
converge. Wittgenstein's insight---that meaning is use, that language games
are forms of life, that understanding requires shared context---is
formalized in the payoff structure. Game theory's insight---that strategic
interaction has equilibria, that cooperation and competition coexist, that
order and information matter---is formalized in the equilibrium analysis.

Neither Wittgenstein nor Nash could have developed this theory alone.
Wittgenstein lacked the mathematical tools; Nash lacked the philosophical
framework. IDT provides both: a mathematics rich enough to capture meaning
creation (through emergence) and a philosophy precise enough to be proved
(through the ideatic space axioms).

The result is a theory of strategic communication that is simultaneously
a philosophy of language (answering Wittgenstein's questions about meaning,
rules, and forms of life), a contribution to game theory (introducing
emergence as a fundamentally new payoff structure), and a practical framework
for analyzing real-world communication (debates, therapy, journalism,
propaganda, courtroom argument, peer review, religious rhetoric).
\end{interpretation}


% ================================================================
% NEW SECTION: MECHANISM DESIGN IN MEANING GAMES
% ================================================================
\section{Mechanism Design: Aggregating Ideas Fairly}

The mechanism design program of Hurwicz (1960), Myerson (1981), and
Maskin (1999) asks: \emph{given a desired social outcome, can we design
a game whose equilibrium achieves it?} In meaning games, the ``social
outcome'' is an aggregated idea, and the ``mechanism'' is a composition
rule.

\begin{definition}[Mechanism Outcome]
\label{def:mechanism-outcome}
A \emph{mechanism} in the meaning game assigns to reports $(a, b)$ the
\emph{mechanism outcome}:
\[
M(a, b) := a \compose b.
\]
This is the simplest aggregation: compose the reports in order. The VCG
principle identifies the ``payment'' each agent must make.
\end{definition}

\begin{definition}[VCG Payment]
\label{def:vcg-payment}
The \emph{VCG payment} for agent 1 (who reports $a$) in the mechanism is:
\[
\text{VCG}_1(a, b) := w(M(a, b)) - w(a) - w(b).
\]
This is the \emph{externality} agent 1 imposes: the difference between the
mechanism outcome's weight and the sum of individual weights.
\end{definition}

\begin{theorem}[VCG Payment Equals Cooperation Surplus]
\label{thm:vcg-coop}
$\text{VCG}_1(a, b) = \Sigma_{\text{coop}}(a, b)$, where the cooperation
surplus is $\Sigma_{\text{coop}}(a, b) := w(a \compose b) - w(a) - w(b)$.
\end{theorem}

\begin{proof}
By direct substitution. The VCG payment formula becomes
$w(a \compose b) - w(a) - w(b)$, which is exactly the cooperation surplus.

The mathematical insight is that in a meaning game, the ``externality''
is the \emph{synergy}: each agent's payment reflects not the harm they
cause others (as in classical VCG) but the \emph{value they co-create}
with others. This is a fundamental difference from standard mechanism
design: in a market for goods, externalities are typically negative
(congestion, pollution); in a meaning game, externalities are typically
positive (emergence, enrichment). The VCG mechanism for meaning games is
not a tax but a \emph{dividend from cooperation}.
\end{proof}

\begin{lstlisting}[caption={Lean 4: VCG mechanism design}]
noncomputable def vcgPayment (report₁ report₂ : I) : ℝ :=
  rs (compose report₁ report₂) (compose report₁ report₂)
  - rs report₂ report₂ - rs report₁ report₁

theorem vcgPayment_eq_coopSurplus (a b : I) :
    vcgPayment a b = cooperationSurplus a b := by
  unfold vcgPayment cooperationSurplus; ring

theorem vcgPayment_void_left (b : I) :
    vcgPayment (void : I) b = 0 := by
  rw [vcgPayment_eq_coopSurplus]; exact cooperationSurplus_void_left b
\end{lstlisting}

\begin{interpretation}[The Hurwicz--Myerson--Maskin Program]
The Nobel-winning mechanism design program asks three questions:
\begin{enumerate}
\item \textbf{Implementation (Maskin):} Can we design a game that
implements a given social choice function in Nash equilibrium?
In IDT, the answer depends on the emergence structure: mechanisms
are implementable in the transparent sector ($\emerge = 0$) and
become increasingly difficult as emergence grows.
\item \textbf{Revenue (Myerson):} How much revenue can an optimal
mechanism extract? In IDT, the ``revenue'' is the cooperation surplus,
which is bounded by the emergence bound (E3):
$|\Sigma_{\text{coop}}| \leq w(a \compose b) \leq w(a) + w(b) + 2\sqrt{w(a) \cdot w(b)}$.
\item \textbf{Incentive compatibility (Hurwicz):} Can agents be
incentivized to report truthfully? The IC gap
$\text{IC}(a, a', b) := u_S^{\text{net}}(a', b) - u_S^{\text{net}}(a, b)$
must be non-positive for all deviations $a'$. At Nash equilibrium,
this is automatically satisfied (Lean theorem T221).
\end{enumerate}

\textbf{What emergence adds:} Classical mechanism design assumes
separable valuations---each agent's value for the outcome depends only
on their own type. In IDT, valuations are \emph{non-separable}:
agent 1's value for the outcome $a \compose b$ depends on the
\emph{interaction} between $a$ and $b$ through the emergence term.
This non-separability is why the attribution impossibility
(Theorem~\ref{thm:attribution-impossibility}) is relevant: the
mechanism designer cannot separately ``price'' each agent's
contribution because the contributions are entangled through emergence.
\end{interpretation}


\section{Auction Theory: Bidding for Ideas}

An \emph{idea auction} is a competitive mechanism where agents bid for
the right to compose with a target idea.

\begin{definition}[Auction Value]
\label{def:auction-value}
The \emph{auction value} for bidder $i$ composing with target $t$ is:
\[
V_i(t) := w(i \compose t) - w(i) = \text{(information content of } t \text{ for } i).
\]
By E4, $V_i(t) \geq 0$: every bidder weakly benefits from composing with
any target.
\end{definition}

\begin{theorem}[Vickrey Truthfulness]
\label{thm:vickrey}
In a second-price (Vickrey) auction for ideas, where the winner pays the
second-highest bid, the \emph{Vickrey payoff} for the winner is:
\[
\pi_V(i, t, p) := V_i(t) - p,
\]
where $p$ is the second-highest bid. This payoff is non-negative whenever
$p \leq V_i(t)$: truthful bidding (bidding one's true value) is
individually rational.
\end{theorem}

\begin{proof}
If $p \leq V_i(t)$, then $\pi_V = V_i(t) - p \geq 0$. The key
economic insight (Vickrey 1961) is that in a second-price auction,
bidding one's true value is a \emph{dominant strategy}: overbidding
risks winning at a loss, and underbidding risks losing a profitable
opportunity. In IDT, the ``true value'' $V_i(t) = w(i \compose t) - w(i)$
is determined by the ideatic space structure, not by subjective preference.
This makes the Vickrey mechanism particularly natural for idea auctions:
the ``true value'' of composing with an idea is an objective quantity
(the information content), so truthful bidding is not just incentive-compatible
but epistemically warranted.
\end{proof}

\begin{lstlisting}[caption={Lean 4: Vickrey auction for ideas}]
/-- Value of composing bidder's idea with target -/
noncomputable def auctionValue (bidder target : I) : ℝ :=
  informationContent target bidder

/-- Auction value is non-negative -/
theorem auctionValue_nonneg (bidder target : I) :
    auctionValue bidder target ≥ 0 :=
  informationContent_nonneg target bidder

/-- Vickrey payoff is non-negative when price ≤ value -/
theorem vickreyPayoff_nonneg (bidder target : I) (price : ℝ)
    (h : price ≤ auctionValue bidder target) :
    auctionValue bidder target - price ≥ 0 := by linarith
\end{lstlisting}

\begin{interpretation}[Market Design and Roth's Repugnance]
Alvin Roth's work on market design emphasizes that some transactions are
considered ``repugnant'' even when they are Pareto-improving (kidney sales,
for example). In IDT, repugnance has a formal correlate: a transaction
$(a, b)$ is ``repugnant'' when the cooperation surplus is positive
($\Sigma_{\text{coop}} > 0$) but the \emph{emergence} is negative for
some important observer $c$: $\emerge(a, b, c) < 0$. The transaction
creates value (by the surplus) but destroys meaning (by the negative
emergence). Roth's insight is that markets require not just efficiency
but \emph{legitimacy}---and legitimacy, in IDT, is positive emergence
with the community's normative context.
\end{interpretation}


\section{Matching Theory: Stable Assignment of Ideas}

Gale and Shapley's (1962) deferred acceptance algorithm finds stable
matchings in two-sided markets. IDT provides a natural quality measure
for matchings: the \emph{match quality} is the cooperation surplus.

\begin{definition}[Match Quality]
\label{def:match-quality}
The \emph{match quality} of pairing idea $a$ with idea $b$ is:
\[
Q(a, b) := \Sigma_{\text{coop}}(a, b) = w(a \compose b) - w(a) - w(b).
\]
A matching is \emph{blocked} by the pair $(a, b)$ if both would prefer
to be matched with each other: $Q(a, b) > Q(a, a')$ and $Q(a, b) > Q(b', b)$
for their current matches $a', b'$.
\end{definition}

\begin{theorem}[Stability and Emergence]
\label{thm:matching-stability}
A matching is \emph{stable} (unblocked) if and only if for every
potential blocking pair $(a, b)$, the cooperation surplus of the
current matching is at least as large as that of the blocking pair.
\end{theorem}

\begin{proof}
By definition: $(a, b)$ blocks the matching if both agents prefer
$(a, b)$ to their current assignment. In IDT, ``prefer'' means
higher cooperation surplus. A matching is stable iff no pair can
achieve higher surplus by re-matching.

The emergence structure adds richness: two ideas might have high
pairwise resonance ($\rs(a, b) > 0$) but low cooperation surplus
(because the emergence $\emerge(a, b, a \compose b) < 0$ is negative).
Such pairs are ``attractive but incompatible''---they resonate but
do not compose well. Conversely, pairs with low pairwise resonance
but high emergence are ``complementary despite dissimilarity.''
The Gale--Shapley algorithm, adapted to IDT, finds the matching that
maximizes total cooperation surplus, which includes both resonance
and emergence.
\end{proof}

\begin{lstlisting}[caption={Lean 4: Matching stability}]
/-- A matching is blocked by pair (a, b) if they prefer each other -/
def isBlocking (a b : I) (currentA currentB : I) : Prop :=
  cooperationSurplus a b > cooperationSurplus a currentA ∧
  cooperationSurplus a b > cooperationSurplus currentB b

/-- An agent cannot block their own self-matching -/
theorem self_stable (a : I) : ¬isBlocking a a a a := by
  unfold isBlocking; intro ⟨h, _⟩; linarith
\end{lstlisting}

\begin{interpretation}[Matching, Marriage, and Intellectual Collaboration]
The matching theory of IDT applies wherever two-sided assignment problems
arise:

\textbf{(a) Academic collaboration.} Researchers seek collaborators who
maximize cooperation surplus---not those who are ``most similar''
(high $\rs(a, b)$) but those who generate the most \emph{emergence}.
The best collaborator is often the one whose expertise is complementary
(moderate $\rs$, high $\emerge$) rather than overlapping (high $\rs$,
low $\emerge$). This explains the empirical finding that
interdisciplinary collaborations, though harder to initiate, often produce
higher-impact work: the emergence from crossing disciplinary boundaries
exceeds the emergence from within-discipline work.

\textbf{(b) Political coalitions.} Parties form coalitions that maximize
political ``weight'' (coalition value). A stable coalition is one that
no pair of parties prefers to defect from. The matching stability theorem
shows that stable coalitions maximize cooperation surplus, which in
political terms means maximizing the emergent policy that exceeds what
each party could achieve alone.

\textbf{(c) The new insight} is that matching quality depends on
\emph{emergence}, not just preference alignment. Two agents with
identical preferences ($\rs(a, b) \gg 0$) but zero emergence
($\emerge(a, b, c) = 0$) produce a match that is ``comfortable but
sterile''---they agree on everything but create nothing new together.
True complementarity requires nonzero emergence: the partner whose ideas,
composed with yours, generate meaning that neither of you could create
alone.
\end{interpretation}

\begin{example}[University Admissions as Matching]
A university ($u$) and a student ($s$) form a ``match'' with quality
$Q(u, s) = w(u \compose s) - w(u) - w(s)$. The admission process is
a matching market. The cooperation surplus captures the ``fit'' between
student and institution: it measures how much the student gains from the
university's intellectual environment \emph{beyond} what they already
know, and vice versa. A high-surplus match is one where the student
grows maximally (high enrichment) \emph{and} the university's
intellectual community is enriched by the student's contributions (high
emergence). This is why ``fit'' matters in admissions: it is not about
test scores ($w(s)$) alone but about the cooperation surplus---the
emergent intellectual value of the student-institution composition.
\end{example}


\section{Network Games: Ideas in Interaction}

When ideas form networks---each idea composed with multiple
partners---the structure of the network determines the emergent
properties of the system.

\begin{definition}[Network Position]
\label{def:network-position}
An agent's \emph{network position} after composing with partners
$p_1, \ldots, p_n$ is:
\[
\text{pos}(a; p_1, \ldots, p_n) := (\cdots((a \compose p_1) \compose p_2)
\cdots) \compose p_n.
\]
The network \emph{value} for agent $a$ is the total enrichment:
$V_{\text{net}}(a) := w(\text{pos}(a; p_1, \ldots, p_n)) - w(a)$.
\end{definition}

\begin{theorem}[Network Value is Non-Negative and Monotone]
\label{thm:network-monotone}
$V_{\text{net}}(a) \geq 0$ for any network, and adding a new partner
$p_{n+1}$ weakly increases the network value:
$V_{\text{net}}(a; p_1, \ldots, p_{n+1}) \geq V_{\text{net}}(a; p_1, \ldots, p_n)$.
\end{theorem}

\begin{proof}
Non-negativity follows from iterated application of E4: each composition
step can only increase weight. For monotonicity, the $(n+1)$-th composition
adds $w(\text{pos}_{n+1}) - w(\text{pos}_n) \geq 0$ by E4. The network
value is a telescoping sum of non-negative increments.

The key insight is that each new network link contributes a
\emph{network externality}: the marginal enrichment from adding partner
$p_{n+1}$ depends on the existing network position $\text{pos}_n$, not
just on $a$ and $p_{n+1}$ individually. This externality is precisely
the emergence $\emerge(\text{pos}_n, p_{n+1}, \text{pos}_{n+1})$. In a
transparent network (zero emergence), externalities are additive---each
link's contribution is independent of other links. In a non-transparent
network, links exhibit \emph{complementarities} (positive emergence:
link $j$ makes link $k$ more valuable) or \emph{substitutabilities}
(negative emergence: link $j$ makes link $k$ less valuable).
\end{proof}

\begin{lstlisting}[caption={Lean 4: Network enrichment}]
/-- Network position: compose agent with all partners -/
def networkPosition (agent : I) : List I → I
  | [] => agent
  | p :: rest => compose (networkPosition agent rest) p

/-- Network enrichment is non-negative -/
theorem network_enrichment (agent partner : I) (rest : List I) :
    rs (networkPosition agent (partner :: rest))
       (networkPosition agent (partner :: rest)) ≥
    rs (networkPosition agent rest)
       (networkPosition agent rest) :=
  compose_enriches (networkPosition agent rest) partner

/-- Network externality is non-negative -/
theorem networkExternality_nonneg (agent : I)
    (existing : List I) (newPartner : I) :
    rs (networkPosition agent (newPartner :: existing))
       (networkPosition agent (newPartner :: existing))
    - rs (networkPosition agent existing)
         (networkPosition agent existing) ≥ 0 := by
  linarith [network_enrichment agent newPartner existing]
\end{lstlisting}

\begin{interpretation}[Networks, Platforms, and the Attention Economy]
The network monotonicity theorem formalizes the \emph{network effect}:
each new connection adds value, and this value can only be positive.
This explains the growth dynamics of social media platforms, academic
citation networks, and cultural memes: participation is always
individually rational (the herding premium is non-negative), so
networks grow until every potential member has joined.

But the weight ratchet adds a darker dimension: the network value is
\emph{irreversible}. Once an agent has composed with the network, they
cannot return to their pre-network state. In the attention economy,
this manifests as the impossibility of ``unplugging'': the ideas,
framings, and emergence patterns absorbed from the network are
permanently part of one's cognitive state. The network does not merely
\emph{inform}---it irreversibly \emph{constitutes}.
\end{interpretation}


\section{Coalition Theory: Ordered Enrichment}

The Deep Games formalization develops a full theory of coalitions where
\emph{order matters}: the sequence in which ideas join the coalition
affects the outcome. This is a fundamental departure from classical
cooperative game theory, where the characteristic function depends
only on the \emph{set} of coalition members.

\begin{definition}[Ordered Coalition]
\label{def:ordered-coalition}
An \emph{ordered coalition} is a list $[a_1, \ldots, a_n]$ of ideas.
Its \emph{coalition value} is the weight of the sequential composition:
\[
V([a_1, \ldots, a_n]) := w(a_1 \compose (a_2 \compose (\cdots \compose a_n) \cdots)).
\]
\end{definition}

\begin{theorem}[Coalition Enrichment]
\label{thm:coalition-enrichment}
For any idea $a$ and any coalition $[a_1, \ldots, a_n]$:
\[
V([a, a_1, \ldots, a_n]) \geq V([a_1, \ldots, a_n]).
\]
Adding a member to the front of a coalition can only increase its value.
\end{theorem}

\begin{proof}
By E4: $w(a \compose C) \geq w(a)$ for any $C$, and symmetrically
$w(a \compose C) \geq w(C)$. The coalition value of $[a, a_1, \ldots, a_n]$
is $w(a \compose C)$ where $C = a_1 \compose (\cdots \compose a_n)$, and
by E4 this is at least $w(C) = V([a_1, \ldots, a_n])$.
\end{proof}

\begin{theorem}[Reorder Effect]
\label{thm:reorder}
The effect of swapping two ideas in a coalition is governed by the
order asymmetry:
\[
\rs(a \compose b, c) - \rs(b \compose a, c) = \emerge(a, b, c) - \emerge(b, a, c).
\]
The reorder effect vanishes if and only if forward and reverse emergence
are equal.
\end{theorem}

\begin{proof}
By the fundamental decomposition:
$\rs(a \compose b, c) = \rs(a, c) + \rs(b, c) + \emerge(a, b, c)$ and
$\rs(b \compose a, c) = \rs(b, c) + \rs(a, c) + \emerge(b, a, c)$.
Subtracting: $\rs(a \compose b, c) - \rs(b \compose a, c)
= \emerge(a, b, c) - \emerge(b, a, c)$.

The result says that the \emph{only} source of order-dependence is the
\emph{asymmetry of emergence}. Two ideas that generate symmetric emergence
($\emerge(a, b, c) = \emerge(b, a, c)$ for all $c$) can be freely reordered
in any coalition without affecting its resonance with any observer.
\end{proof}

\begin{lstlisting}[caption={Lean 4: Coalition theory}]
/-- Ordered coalition: compose all ideas in sequence -/
def orderedCoalition : List I → I
  | [] => void
  | a :: rest => compose a (orderedCoalition rest)

/-- Coalition weight is at least the weight of the first member -/
theorem coalWeight_ge_first' (a : I) (rest : List I) :
    rs (orderedCoalition (a :: rest)) (orderedCoalition (a :: rest)) ≥
    rs a a :=
  compose_enriches' a (orderedCoalition rest)

/-- Reorder effect is the emergence asymmetry -/
theorem reorder_effect (a b obs : I) :
    rs (compose a b) obs - rs (compose b a) obs =
    emergence a b obs - emergence b a obs := by
  unfold emergence; ring
\end{lstlisting}

\begin{remark}[The Paradox of Collective Authorship]
Who wrote the coalition's output? The reorder theorem shows that
attribution depends on order: if $a$ joins before $b$, the emergence
$\emerge(a, b, c)$ is attributed differently than if $b$ joins first
($\emerge(b, a, c)$). This connects to the attribution impossibility
(Theorem~\ref{thm:attribution-impossibility}): when emergence is
nonzero, there is no consistent way to assign credit to individual
members. Every collaborative work---from a jointly authored paper to
a legislative compromise to a jazz improvisation---faces this paradox.
The solution is not to attribute more precisely but to recognize that
emergence is irreducibly collective.
\end{remark}

\begin{example}[The Founding Fathers: Order of Constitutional Arguments]
At the Constitutional Convention of 1787, the order in which proposals
were introduced shaped the final document. Madison's Virginia Plan,
presented first, set the frame ($a_1$). Paterson's New Jersey Plan
($a_2$), composed with Madison's frame, produced different emergence
than if Paterson had spoken first. The Great Compromise was not merely
the set $\{a_1, a_2\}$ but the \emph{ordered} coalition $[a_1, a_2]$,
and the reorder theorem tells us that $V([a_1, a_2]) \neq V([a_2, a_1])$
unless the two plans generate symmetric emergence. History---the order
in which ideas are proposed---is mathematically consequential.
\end{example}

\begin{remark}[Connection to Category Theory: Cocycle Algebras]
The ordered coalition structure, combined with the cocycle condition
on emergence (Theorem~\ref{thm:cocycle}), gives the set of ideas the
structure of a \emph{cocycle algebra}. Formally, the emergence function
$\emerge : \IS \times \IS \times \IS \to \R$ satisfies a cocycle-like
identity:
\[
\emerge(a \compose b, c, d) = \emerge(a, c, d) + \emerge(b, c, d)
+ \text{higher-order terms}.
\]
This is analogous to the group cocycle condition in algebraic topology,
where cocycles classify extensions of groups. In IDT, the emergence
cocycle classifies the ``extensions'' of meaning: how the combination
of ideas produces meaning beyond the sum of parts. The vanishing of
the cocycle (transparency) corresponds to a trivial extension; nonzero
emergence corresponds to a nontrivial extension of the meaning algebra.
\end{remark}

\begin{interpretation}[Coalition Theory and Political Science]
The ordered coalition theory has direct applications to political science:

\textbf{(a) Legislative bargaining:} In Romer--Rosenthal (1978) agenda
models, the order in which bills are considered affects the outcome.
IDT's reorder theorem makes this precise: the ``agenda effect'' is
exactly the emergence asymmetry $\emerge(a, b, c) - \emerge(b, a, c)$.
A powerful agenda-setter chooses the order that maximizes their
preferred emergence pattern.

\textbf{(b) Coalition formation:} In Riker's (1962) theory of political
coalitions, the ``minimum winning coalition'' is the smallest group that
can achieve the outcome. IDT adds that coalition \emph{order} matters:
the same parties, joining in different orders, produce different emergent
policies. This explains why the sequence of campaign endorsements matters
more than the final list of endorsers.

\textbf{(c) Deliberative democracy:} Habermas's ideal of rational
deliberation requires that the order of speakers not affect the outcome.
The reorder theorem shows that this requirement is equivalent to
demanding symmetric emergence: $\emerge(a, b, c) = \emerge(b, a, c)$
for all participants. In practice, this is rarely satisfied---which is
why procedural rules (Robert's Rules of Order, parliamentary procedure)
exist to \emph{mitigate} but cannot \emph{eliminate} order effects.
\end{interpretation}

\begin{interpretation}[The Synthesis: From Wittgenstein to Market Design]
The seven chapters of this volume trace a remarkable arc:

\textbf{Chapter 1} formalized Wittgenstein's philosophy of language
in IDT, showing that ``meaning is use'' corresponds to resonance profiles,
``family resemblance'' to non-transitive resonance chains, and ``forms
of life'' to shared contexts with positive resonance.

\textbf{Chapters 2--3} developed the game-theoretic structure of meaning,
showing that Nash equilibria exist but have novel properties
(emergence-dependent, order-sensitive, potentially babbling).

\textbf{Chapter 4} proved that the communication surplus IS emergence,
and that the attribution of this surplus is fundamentally impossible
when emergence is nonzero---connecting to Shapley's cooperative game
theory and revealing its limitations.

\textbf{Chapter 5} analyzed persuasion through the weight ratchet,
showing that repeated exposure creates irreversible cognitive change,
and that information cascades emerge from the non-negativity of herding
premiums.

\textbf{Chapter 6} formalized deception as negative-emergence engineering
and strategic silence as a theorem, connecting to Grice's maxims and
Eco's theory of interpretation.

\textbf{Chapter 7} extended the framework to two-sided games, mechanism
design, auctions, matching, networks, and coalitions---showing that the
emergence term transforms every classical result in non-trivial ways.

The unifying insight is this: \textbf{emergence is the mathematics of
meaning}. Wherever meaning is created---in conversation, persuasion,
deception, markets, networks, coalitions---the emergence term
$\emerge(a, b, c)$ captures exactly what classical theories miss.
Von Neumann's game theory is the $\emerge = 0$ special case.
Shannon's information theory is the $\emerge = 0$ special case.
Shapley's cooperative game theory is the $\emerge = 0$ special case.
The non-trivial case---the case that matters for every real human
interaction---is $\emerge \neq 0$. This is the territory that IDT maps.
\end{interpretation}

\begin{remark}[Looking Forward: Open Questions]
This volume has established the game-theoretic foundations of meaning. But
several deep questions remain open:

\begin{enumerate}
\item \textbf{Computational complexity.} Given an ideatic space with $n$
ideas, how hard is it to find the optimal signal in a meaning game? The
general problem is likely NP-hard (the Lean formalization proves a related
NP-hardness result for learning in ideatic spaces). But special cases---
transparent games, games with bounded emergence---may be tractable.

\item \textbf{Evolutionary stability.} The evolutionary dynamics section
showed that fitness equals welfare, but the question of which strategies
are \emph{evolutionarily stable} (Maynard Smith 1982) in meaning games
remains open. The weight ratchet suggests that ESS analysis may differ
fundamentally from the classical case: once an idea has been ``learned''
(composed into the population's state), it cannot be ``unlearned.''

\item \textbf{Large-scale meaning games.} What happens when thousands or
millions of agents play meaning games simultaneously? The network and
coalition results provide initial answers, but a full theory of
\emph{meaning markets} (where ideas are traded, priced, and valued) is
still to be developed. Volume~V (``The Mathematics of Power'') takes up
this challenge.

\item \textbf{The dynamics of emergence.} We have treated $\emerge$ as
a static function, but in real communities, the emergence landscape
evolves as the community's shared context changes. A dynamical theory
of emergence---tracking how $\emerge(a, b, c)$ changes as $b$ evolves
under repeated meaning games---would connect IDT to the theory of
cultural evolution and historical change.
\end{enumerate}
\end{remark}
